\article{线性群的主 L 函数}{Principal L-functions of the linear group}{Hervé Jacquet}{2}{63--86}
设 $F$ 为局部域或整体域（local or global field），将 $G_n=\GL(n)$ 视为 $F$ 上的代数群。此时 ${}^LG_n^0=\GL(n,\C)$。记后者在 $\C^n$ 上的自然表示（natural representation）为 $r_n$；与 $r_n$ 对应的 $L$ 函数在“Langlands 哲学”（Langlands philosophy）中居于核心地位。本文结合近期关于表示分类的结果，回顾这类 $L$ 函数的部分理论。

\origsection{1}{非 Archimedes 局部理论（Local nonarchimedean theory）}
以下设 $F$ 为非 Archimedes 局部域。在不致混淆时，简记 $G_n(F)$ 为 $G_n$；对其他 $F$-群也采用这一约定。

\paragraph*{(1.1)} 设 $\pi$ 为 $G_n$ 在复向量空间 $V$ 上的容许表示（admissible representation）；以 $\widetilde\pi$ 表示其逆反表示（contragredient representation），$\widetilde V$ 为后者作用的空间，$\langle\cdot,\cdot\rangle$ 为 $V\times\widetilde V$ 上的典范不变双线性形式。$\widetilde\pi$ 容许，且 $\pi$ 不可约时它也不可约。此外 $\widetilde{\widetilde\pi}\simeq\pi$。函数
\begin{equation}
 g\longmapsto\langle\pi(g)v,\widetilde v\rangle\qquad(v\in V,\widetilde v\in\widetilde V)\tag{1.1.1}
\end{equation}
及其线性组合称为 $\pi$ 的矩阵系数（coefficients），下文简称系数。若 $f$ 为 $\pi$ 的系数，则由
\begin{equation}
 f^\vee(g)=f(g^{-1})\tag{1.1.2}
\end{equation}
定义的函数为 $\widetilde\pi$ 的系数。

令 $M(p\times q,F)$ 为元素在 $F$ 中的 $p$ 行 $q$ 列矩阵空间；以 $\mathcal S(p\times q,F)$ 表示该空间上的 Schwartz--Bruhat 函数（Schwartz--Bruhat functions）空间。

若 $f$ 为 $\pi$ 的系数，$\Phi\in\mathcal S(n\times n,F)$，$s\in\C$，令
\begin{equation}
 Z(\Phi,s,f)=\int\Phi(g)|\det g|^s f(g)\,d^\times g.\tag{1.1.3}
\end{equation}
积分遍历 $G_n$，$d^\times g$ 为该群上的 Haar 测度。下面重述 [R.G.--H.J.] 的结果（命题(1.2)、(1.4)）。

\begin{originaltheorem}{命题 1.2（Proposition 1.2）}
假定 $\pi$ 不可约。

(1) 存在 $s_0$，使积分(1.1.3)在半平面 $\operatorname{Re}s>s_0$ 绝对收敛。

(2) 若 $F$ 的剩余域有 $q$ 个元素，则这些积分给出 $q^{-s}$ 的有理函数，且具有不依赖 $f,\Phi$ 的公分母。

(3) 设 $\psi\neq1$ 为 $F$ 的加性特征（additive character）。存在有理函数 $\gamma(s,\pi,\psi)$，使对 $\pi$ 的所有系数 $f$ 和所有 $\Phi$，有
\begin{equation}
 Z(\widehat\Phi,1-s+\tfrac12(n-1),f^\vee)
 =\gamma(s,\pi,\psi)Z(\Phi,s,f),\tag{1.2.3}
\end{equation}
其中 $\widehat\Phi$ 为 $\Phi$ 关于 $\psi$ 的 Fourier 变换（Fourier transform）。\footnote{\textbf{译者注（本次校订）：}按 (1.1.3) 对 $Z$ 的定义及 (1.3.7) 的正规化，本式右端应为 $\gamma(s,\pi,\psi)Z(\Phi,s+\tfrac12(n-1),f)$。两边积分的参数分别为 $1-s+\tfrac12(n-1)$ 与 $s+\tfrac12(n-1)$，其和为 $n$；这里的右端漏了平移项。第 4 节的 (4.2.3) 有同一问题。}
\end{originaltheorem}
在(3)中，$\widehat\Phi$ 由
\begin{equation}
 \widehat\Phi(x)=\int\Phi(y)\psi[\operatorname{tr}(yx)]\,dy\tag{1.2.4}
\end{equation}
定义，其中 $dy$ 为 $M(n\times n,F)$ 上的自对偶（self-dual）Haar 测度。将 (1)、(2) 应用于 $\widetilde\pi$，便可解释 (3) 左端的亚纯延拓。命题 (1.2) 也可以直接对二元组 $(\pi,\widetilde\pi)$ 表述：它对 $\pi$ 与 $\widetilde\pi$ 是对称的，证明的每一步都保留这一对称性。

\paragraph*{例 1.2.5（Example 1.2.5）。}
设 $\pi$ 为尖点表示（cuspidal）。当 $n=1$ 时，尖点条件自动成立；此时 $\pi$ 是 $F^\times$ 的拟特征（quasi-character），命题 (1.2) 即为 [J.T.1] 或 [A.W.3] 中的结果。当 $n>1$ 时，$\pi$ 的系数模去 $G_n$ 的中心后具有紧支集。利用这一性质即可证明命题 (1.2)，方法与 $n=1$ 时基本相同（[R.G.--H.J.]、[H.J.1]）。

\paragraph*{(1.3)} 命题 (1.2) 的证明留待第 2 节，这里先推导几个简单结论。若 $f,\Phi$ 如上，$h\in G_n$，则由
\begin{equation}
 f_1(g)=f(gh),\qquad\Phi_1(x)=\Phi(xh)\tag{1.3.1}
\end{equation}
定义的函数仍为同一类型。且若假定(1.2.1)、(1.2.2)，则
\begin{equation}
 Z(\Phi_1,s,f_1)=|\det h|^{-s}Z(\Phi,s,f).\tag{1.3.2}
\end{equation}
因此，由积分
\begin{equation}
 Z(\Phi,s+\tfrac12(n-1),f)\tag{1.3.3}
\end{equation}
张成的 $\C(q^{-s})$ 的向量子空间 $I(\pi)$ 实际上是环 $\C[q^{-s},q^s]$ 的分式理想（fractional ideal）。取 $f(e)\neq0$，再适当选取支集位于 $e$ 的充分小邻域内的 $\Phi$，便可使 (1.3.3) 成为与 $s$ 无关的非零常数。因此 $I(\pi)$ 包含所有常数，从而有形如 $P(q^{-s})^{-1}$ 的生成元，其中 $P\in\C[X]$。要求 $P(0)=1$ 加以正规化，并令
\begin{equation}
 L(s,\pi)=P(q^{-s})^{-1}.\tag{1.3.4}
\end{equation}
\paragraph*{例 1.3.5。}
若 $\pi$ 超尖点（supercuspidal），则 $L(s,\pi)=1$，除非 $n=1$ 且 $\pi(x)=|x|^t$。此时（同上）$L(s,\pi)=(1-q^{-s-t})^{-1}$。

现假定命题(1.2)对 $(\pi,\widetilde\pi)$ 成立。因 $I(\pi)\neq0$，因子 $\gamma$ 唯一（这也使该记号合理）。令
\begin{equation}
 \varepsilon(s,\pi,\psi)=\gamma(s,\pi,\psi)L(s,\pi)/L(1-s,\widetilde\pi).\tag{1.3.6}
\end{equation}
于是(1.2.3)写为
\begin{equation}
 \frac{Z(\widehat\Phi,1-s+\tfrac12(n-1),f^\vee)}{L(1-s,\widetilde\pi)}
 =\varepsilon(s,\pi,\psi)\frac{Z(\Phi,s+\tfrac12(n-1),f)}{L(s,\pi)}.\tag{1.3.7}
\end{equation}
由 $L$ 因子的定义，$\varepsilon(s,\pi,\psi)$ 及其逆均在 $\C[q^{-s},q^s]$ 中，故 $\varepsilon(s,\pi,\psi)$ 是 $q^{-s}$ 的单项式。

令 $\omega$ 为 $\pi$ 的中心拟特征，即满足下式的同态 $\omega:F^\times\to\C^\times$：
\begin{equation}
 \pi(a)=\pi(a1_n)=\omega(a)1_V\qquad(a\in F^\times).\tag{1.3.8}
\end{equation}
若(1.2.3)对一种 $\psi$ 成立，则对其他任意 $\psi'$ 也成立。特别地，
\begin{equation}
 \varepsilon(s,\pi,\psi')=\omega(b)|b|^{(s-1/2)n}\varepsilon(s,\pi,\psi)
 \quad\text{若 }\psi'(x)=\psi(xb).\tag{1.3.9}
\end{equation}
交换 $\pi,\widetilde\pi$，还得到（若命题(1.2)对该二元组成立）
\begin{equation}
 \varepsilon(1-s,\widetilde\pi,\psi)\varepsilon(s,\pi,\psi)=\omega(-1).\tag{1.3.10}
\end{equation}
此外，对任意拟特征 $\chi$，以 $\pi\otimes\chi$ 表示 $G_n$ 在 $V$ 上的表示 $g\mapsto\pi(g)\chi(\det g)$。以 $\alpha_F$ 或 $\alpha$ 表示 $F$ 的模（module）。则
\begin{align}
 L(s,\pi\otimes\alpha^t)&=L(s+t,\pi),\tag{1.3.11}\\
 \varepsilon(s,\pi\otimes\alpha^t,\psi)&=\varepsilon(s+t,\pi,\psi).\tag{1.3.12}
\end{align}
$\gamma$ 因子满足类似恒等式。

\origsection{2}{诱导表示（Induced representations）}
$F$ 仍为非 Archimedes 局部域。

\paragraph*{(2.1)} 令 $P$ 为 $G_n$ 的 $F$-抛物子群。具体取类型为 $(n_1,n_2,\ldots,n_r)$ 的标准抛物子群，其矩阵 $p$ 形如
\begin{equation}
 p=\begin{pmatrix}m_1&*&\cdots&u_{ij}\\0&m_2&\ddots&\vdots\\\vdots&\ddots&\ddots&*\\0&\cdots&0&m_r\end{pmatrix},\qquad m_i\in G_{n_i}.\tag{2.1.1}
\end{equation}
并令
\begin{equation}
 U=U_P\quad(P\text{ 的幺幂根基（unipotent radical）}),\qquad M=M_P=P/U_P.\tag{2.1.2}
\end{equation}
通常将 $M$ 与 $P$ 中所有非对角分块 $u_{ij}$ 都为零的矩阵所成的子群等同。这样便有自然同构
\begin{equation}
 M\simeq\prod_iG_{n_i}\qquad(1\leq i\leq r).\tag{2.1.3}
\end{equation}
给定各 $G_{n_i}$ 的容许表示 $\sigma_i$，先构造 $M=P/U$ 的表示 $\sigma=\mathop{\times}_i\sigma_i$，再将其拉回为 $P$ 的表示（$U$ 在其中平凡作用），最后诱导到 $G_n$。所得表示
\begin{equation}
 \xi=I(G,P;\sigma)=I(G,P;\sigma_1,\sigma_2,\ldots,\sigma_r)\tag{2.1.4}
\end{equation}
容许。其逆反 $\widetilde\xi$ 等价于
\begin{equation}
 \xi'=I(G,P;\widetilde\sigma)=I(G,P;\widetilde\sigma_1,\widetilde\sigma_2,\ldots,\widetilde\sigma_r).\tag{2.1.5}
\end{equation}
更准确地，令 $W$ 为 $\sigma$ 的空间，则 $\xi$ 的空间 $V$ 由从 $G_n$ 到 $W$ 的所有光滑且满足
\begin{equation}
 F(pg)=\delta_P(p)^{1/2}\sigma(p)F(g)\tag{2.1.6}
\end{equation}
的函数 $F$ 构成，其中 $\delta_P$ 为 $P$ 的模。群 $G_n$ 通过右平移作用于 $V$。$\xi'$ 的空间 $V'$ 由 $\widetilde\sigma$ 的空间 $\widetilde W$ 类似定义。对 $F\in V,\widetilde F\in V'$，函数 $\varphi(g)=\langle F(g),\widetilde F(g)\rangle$ 满足
\begin{equation}
 \varphi(pg)=\delta_P(p)\varphi(g).\tag{2.1.7}
\end{equation}
用
\begin{equation}
 \varphi\longmapsto\int_{P\backslash G}\varphi(g)\,dg\tag{2.1.8}
\end{equation}
表示满足(2.1.7)的连续函数空间上的一个正的右不变形式，便可令
\begin{equation}
 \langle F,\widetilde F\rangle=\int_{P\backslash G}\langle F(g),\widetilde F(g)\rangle\,dg.\tag{2.1.9}
\end{equation}
此配对使我们可将 $V'$ 与 $\widetilde V$、$\xi'$ 与 $\widetilde\xi$ 等同。令 $R_F$ 为 $F$ 的整数环，取
\begin{equation}
 K=K_n=\GL(n,R_F).\tag{2.1.10}
\end{equation}
因 $G_n=PK_n$，可将(2.1.8)取为
\begin{equation}
 \varphi\longmapsto\int_{K_n}\varphi(k)\,dk.\tag{2.1.11}
\end{equation}

\paragraph*{(2.2)} 为便于计算，下面用 $\sigma$ 的系数来描述 $\xi$ 的系数。设 $f$ 为 $F\in V,\widetilde F\in\widetilde V=V'$ 所确定的 $\xi$ 的系数（(1.1.1)）。则由 $H(g_1,g_2)=\langle F(g_1),\widetilde F(g_2)\rangle$ 定义的函数 $H:G_n\times G_n\to\C$ 满足
\begin{align}
 H(u_1mg_1,u_2mg_2)&=\delta_P(m)H(g_1,g_2),\quad
 g_i\in G_n,u_i\in U_P,m\in M_P;\tag{2.2.1}\\
 m\longmapsto H(mg_1,g_2)&\text{ 对任意 }g_1,g_2\text{ 均为 }\sigma\otimes\delta_P^{1/2}\text{ 的系数};\tag{2.2.2}\\
 H&\text{ 在右侧 }K_n\times K_n\text{-有限}.\tag{2.2.3}
\end{align}
此外，
\begin{equation}
 f(g)=\int_{P\backslash G}H(hg,h)\,dh=\int_KH(kg,k)\,dk.\tag{2.2.4}
\end{equation}
反之，若 $H$ 为满足(2.2.1)--(2.2.3)的任意函数，则(2.2.4)定义的 $f$ 为 $\xi$ 的系数。$\widetilde\xi$ 的系数 $f^\vee$ 相应地由
\begin{equation}
 f^\vee(g)=\int_{P\backslash G}\widetilde H(hg,h)\,dh,\qquad
 \widetilde H(g_1,g_2)=H(g_2,g_1)\tag{2.2.5}
\end{equation}
给出，且 $\widetilde H$ 满足将 $\sigma$ 换为 $\widetilde\sigma$ 后的(2.2.1)--(2.2.3)。

\paragraph*{(2.3)} 在陈述本节的主要定理之前，先作一点说明：对于具有中心拟特征（见 (1.3.8)）的容许表示 $\pi$，即使不要求它不可约，仍可对 $(\pi,\widetilde\pi)$ 表述 (1.2) 中的各项断言；当然，此时它们未必成立。以下假定每个 $\sigma_i$ 都有中心拟特征，因此 $\xi$ 也有中心拟特征。

\begin{originaltheorem}{命题 2.3（Proposition 2.3）}
沿用(2.1)的记号，假定(1.2)的断言对每个 $(\sigma_i,\widetilde\sigma_i)$ 成立。则它们对 $(\xi,\widetilde\xi)$ 也成立。此外，
\[
 I(\xi)=\prod_iI(\sigma_i),\qquad
 \gamma(s,\xi,\psi)=\prod_i\gamma(s,\sigma_i,\psi).
\]
\end{originaltheorem}
证明分为 (2.4)--(2.6) 三步。

\paragraph*{(2.4)} 令 $f$ 由(2.2.4)给出。交换积分次序，得
\begin{equation}
\begin{aligned}
 Z(\Phi,s+\tfrac12(n-1),f)
 &=\int_Kdk\int_{G_n}\Phi(g)|\det g|^{s+(n-1)/2}H(kg,k)\,d^\times g\\
 &=\int_Kdk\int_{G_n}\Phi(k^{-1}g)|\det g|^{s+(n-1)/2}H(g,k)\,d^\times g\\
 &=\int dk\,dk'\int_P\Phi(k^{-1}pk')|\det p|^{s+(n-1)/2}H(pk',k)\,d_lp.
\end{aligned}\tag{2.4.1}
\end{equation}
再定义（$p$ 如(2.1.1)）
\begin{align}
 \Psi(m_1,\ldots,m_r;k,k')&=\int\Phi(k^{-1}pk')\bigotimes_{i<j}du_{ij},\tag{2.4.2}\\
 h(m_1,\ldots,m_r;k,k')&=H(pk',k)\delta_P^{-1/2}(p),\tag{2.4.3}
\end{align}
对变量 $u_{ij}$ 积分后，(2.4.1)可写为
\begin{equation}
 \int dk\,dk'\int\Psi(m_1,\ldots,m_r;k,k')h(m_1,\ldots,m_r;k,k')
 \prod_i|\det m_i|^{s+(n_i-1)/2}\bigotimes_i d^\times m_i.\tag{2.4.4}
\end{equation}
由于所涉及函数 $K$-有限，对于给定 $\Phi,H$，(2.4.4)可以写成内层积分在 $K\times K$ 中某个有限集上的和。给定 $k,k'$ 后，$\Psi$ 是若干乘积 $\prod_i\Psi_i(m_i)$ 的有限和，其中 $\Psi_i\in\mathcal S(n_i\times n_i,F)$；同样，$h$ 是若干乘积 $\prod_if_i(m_i)$ 的有限和，其中 $f_i$ 为 $\sigma_i$ 的系数（(2.2.2)）。因此(2.4.1)为若干乘积
\begin{equation}
 \prod_i Z(\Phi_i,s+\tfrac12(n_i-1),f_i)\tag{2.4.5}
\end{equation}
的有限和，其中 $\Phi_i\in\mathcal S(n_i\times n_i,F)$，$f_i$ 为 $\sigma_i$ 的系数。这就证明了 $\xi$ 满足 (1.2.1)、(1.2.2)，并且得到包含关系 $I(\xi)\subset\prod_iI(\sigma_i)$。

\paragraph*{(2.5)} 下面证明反向包含关系。给定形如 (2.4.5) 的表达式，可选取 $\Phi\in\mathcal S(n\times n,F)$，使按(2.1.1)的记号有
\begin{equation}
 \int\Phi(p)\bigotimes_{i<j}du_{ij}=\prod_i\Phi_i(m_i).\tag{2.5.1}
\end{equation}
此外，$K$ 上存在两个 $K$-有限函数 $\eta,\eta'$，使
\begin{equation}
 \iint\Phi(k^{-1}xk')\eta(k)\eta'(k')\,dk\,dk'=\Phi(x).\tag{2.5.2}
\end{equation}
于是当 $\operatorname{Re}s$ 充分大时，(2.4.5)等于
\begin{equation}
 \int\Phi(k^{-1}pk')|\det p|^{s+(n-1)/2}
 \prod_if_i(m_i)\delta_P^{1/2}(p)\eta(k)\eta'(k')\,d_lp\,dk\,dk'.\tag{2.5.3}
\end{equation}
令 $dh$ 为紧群 $K\cap P$ 上正规化 Haar 测度。将 $k$ 换成 $hk$，$k'$ 换成 $h'k'$，其中 $h,h'\in K\cap P$，再对 $(K\cap P)\times(K\cap P)$ 积分，得到(2.5.3)等于
\begin{equation}
\begin{aligned}
 \int dh\,dh'\int&\Phi(k^{-1}h^{-1}ph'k')|\det p|^{s+(n-1)/2}\\
 &\cdot\prod_i f_i(m_i)\delta_P^{1/2}(p)\eta(hk)\eta'(h'k')\,d_lp\,dk\,dk'.
\end{aligned}\tag{2.5.4}
\end{equation}
现将 $p$ 换成 $hph'^{-1}$，并将 $h,h'$ 写为(2.1.1)的形式，以 $h_i\in K_{n_i},h'_i\in K_{n_i}$ 代替 $m_i$。则(2.5.4)等于
\begin{equation}
 \int\Phi(k^{-1}pk')H_1(p,k,k')|\det p|^{s+(n-1)/2}\,d_lp\,dk\,dk',\tag{2.5.5}
\end{equation}
其中
\[
 H_1(p,k,k')=\iint\eta(hk)\eta'(h'k')\prod_i f_i(h_im_ih_i'^{-1})\delta_P^{1/2}(p)\,dh\,dh'.
\]
容易验证，存在满足(2.2.1)--(2.2.3)的函数 $H$，使 $H(pk,k')=H_1(p,k,k')$（$p\in P,k,k'\in K$）。若 $f$ 为对应的 $\xi$ 的系数（(2.2.4)），比较(2.5.4)与(2.4.1)，得到
\[
 \prod_i Z(\Phi_i,s+\tfrac12(n_i-1),f_i)=Z(\Phi,s+\tfrac12(n-1),f).
\]
因此 $I(\xi)=\prod_iI(\sigma_i)$。

\paragraph*{(2.6)} 下面证明函数方程。在(2.4.1)中用 $f^\vee$ 代 $f$，$\widehat\Phi$ 代 $\Phi$。于是 $H$ 换为 $\widetilde H$（(2.2.5)），恒等式(2.4.3)给出
\begin{equation}
 h(m_1^{-1},\ldots,m_r^{-1};k',k)=\widetilde H(pk,k')\delta_P^{-1/2}(p).\tag{2.6.1}
\end{equation}
类似地，(2.4.2)给出
\begin{equation}
 \widehat\Psi(m_1,\ldots,m_r;k',k)=\int\widehat\Phi(k^{-1}pk')\bigotimes_{i<j}du_{ij},\tag{2.6.2}
\end{equation}
其中 $\widehat\Psi$ 表示 $\Psi$ 关于每个变量 $m_i$ 的 Fourier 变换。代替(2.4.1)、(2.4.4)，对 $\operatorname{Re}s$ 充分大，现在有
\begin{equation}
\begin{aligned}
 Z(\widehat\Phi,s+\tfrac12(n-1),f^\vee)
 &=\int dk\,dk'\int\widehat\Phi(k^{-1}pk')|\det p|^{s+(n-1)/2}\widetilde H(pk',k)\,d_lp\\
 &=\int dk\,dk'\int\widehat\Psi(m_1,\ldots,m_r;k',k)h(m_1^{-1},\ldots,m_r^{-1};k',k)\\
 &\hspace{3em}\cdot\prod_i|\det m_i|^{s+(n_i-1)/2}\bigotimes_i d^\times m_i.
\end{aligned}\tag{2.6.3}
\end{equation}
将各表示 $\sigma_i$ 的函数方程代入，并利用上面对 $h,\Psi$ 的分解，便得到
\[
 Z(\widehat\Phi,1-s+\tfrac12(n-1),f^\vee)
 =\prod_i\gamma(s,\sigma_i,\psi)Z(\Phi,s+\tfrac12(n-1),f).
\]
至此完成(2.3)的证明。

\paragraph*{(2.7)} 仍用(2.3)的记号。将(2.3)用于 $\xi,\widetilde\xi$，得到
\begin{equation}
 L(s,\xi)=\prod_iL(s,\sigma_i),\qquad
 L(s,\widetilde\xi)=\prod_iL(s,\widetilde\sigma_i).\tag{2.7.1}
\end{equation}
于是(2.3)最后一个断言写为
\begin{equation}
 \varepsilon(s,\xi,\psi)=\prod_i\varepsilon(s,\sigma_i,\psi).\tag{2.7.2}
\end{equation}
令 $\pi$ 为 $\xi$ 的一个不可约组成成分。则 $\pi$ 的每个系数都是 $\xi$ 的系数，且 $\widetilde\pi$ 为 $\widetilde\xi$ 的不可约组成成分。因此命题(1.2)对 $(\pi,\widetilde\pi)$ 成立。此外，
\begin{equation}
 I(\pi)\subset I(\xi),\quad I(\widetilde\pi)\subset I(\widetilde\xi),\quad
 \gamma(s,\pi,\psi)=\gamma(s,\xi,\psi).\tag{2.7.3}
\end{equation}
故存在两个多项式 $P,\widetilde P\in\C[X]$，使
\begin{equation}
\begin{gathered}
 L(s,\pi)=P(q^{-s})L(s,\xi),\qquad
 L(s,\widetilde\pi)=\widetilde P(q^{-s})L(s,\widetilde\xi),\\
 P(0)=\widetilde P(0)=1.
\end{gathered}\tag{2.7.4}
\end{equation}
注意，$P(q^{-s})$ 在 $\C[q^{-s}]$ 中整除 $L(s,\xi)^{-1}$。此外，
\begin{equation}
\begin{aligned}
 \varepsilon(s,\pi,\psi)
 &=\gamma(s,\pi,\psi)\frac{L(s,\pi)}{L(1-s,\widetilde\pi)}\\
 &=\gamma(s,\xi,\psi)\frac{L(s,\xi)P(q^{-s})}{L(1-s,\widetilde\xi)\widetilde P(q^{-1+s})}\\
 &=\varepsilon(s,\xi,\psi)\frac{P(q^{-s})}{\widetilde P(q^{-1+s})}.
\end{aligned}\tag{2.7.5}
\end{equation}
因 $\varepsilon$ 因子是环 $\C[q^{-s},q^s]$ 的单位（units），有
\begin{equation}
 P(X)=\prod_i(1-a_iX),\qquad
 \widetilde P(X)=\prod_i(1-a_i^{-1}qX).\tag{2.7.6}
\end{equation}

\paragraph*{(2.8)} 一般地，若 $\pi$ 为 $G_n$ 的任意不可约容许表示，它是形如(2.1.4)、其中各 $\sigma_i$ 尖点的诱导表示的组成成分。因(1.2)对每个 $\sigma_i$ 成立（(1.2.5)），故对 $(\xi,\widetilde\xi)$ 成立（(2.3)），从而对 $(\pi,\widetilde\pi)$ 成立。至此，命题 (1.2) 的证明完成。

\origsection{3}{L 因子的计算（Computation of the L-factor）}
对于 $G_n$ 的任意不可约容许表示 $\pi$，命题 (1.2) 已经得到证明，$\gamma(s,\pi,\psi)$ 也已确定（见 (2.7.3) 及命题 2.3）。现在只需计算 $L$ 因子。

\origsubsection{3.1}{平方可积表示（Square-integrable representations）}
已见，若 $\pi$ 尖点，则 $L(s,\pi)=1$，除非 $n=1$、$\pi=\alpha^t$；此时
\begin{equation}
 L(s,\pi)=(1-q^{-s-t})^{-1}.\tag{3.1.1}
\end{equation}
下面计算本质平方可积表示（essentially square-integrable representation）的 $L$ 因子。回顾其定义：若 $\pi$ 有中心特征 $\omega$，且其各个系数（按 $\omega$ 变换）在模中心的商空间上平方可积，就称 $\pi$ 为平方可积表示；若 $\pi=\pi_0\otimes\alpha^t$，其中 $\pi_0$ 平方可积、$t\in\R$，就称 $\pi$ 为本质平方可积表示。先回顾 Bernstein 和 Zelevinsky 对这类表示的构造。

令 $r\mid n$，$n=jr$。令 $P$ 为 $G_n$ 中类型 $(j,j,\ldots,j)$ 的标准抛物子群，$\tau$ 为 $G_j$ 的尖点表示。对 $1\leq i\leq r$ 令 $\sigma_i=\tau\otimes\alpha^{i-1}$。则诱导表示
\begin{equation}
 \xi=I(G,P;\sigma_1,\sigma_2,\ldots,\sigma_r)\tag{3.1.2}
\end{equation}
具有唯一的本质平方可积组成成分 $\pi$。所有本质平方可积表示均如此得到，且 $r,\tau$ 由 $\pi$ 唯一确定。

\begin{originaltheorem}{命题 3.1.3（Proposition 3.1.3）}
沿用上述记号，有 $L(s,\pi)=L(s,\tau)$。\footnote{\textbf{译者注（本次校订）：}这里的公式与前文 $\sigma_i=\tau\otimes\alpha^{i-1}$ 的递增约定不一致。按这一约定，应为 $L(s,\pi)=L(s+r-1,\tau)$；若将构造中的指数改为 $1-i$，才得到正文所写的 $L(s,\tau)$。例如 $r=2$、$\tau=\alpha^{-1/2}$ 时，平方可积成分为酉 Steinberg 表示，其因子为 $(1-q^{-s-1/2})^{-1}$。正文公式却给出 $(1-q^{-s+1/2})^{-1}$，与紧接着的注 3.1.4 所述极点位置相矛盾。}
\end{originaltheorem}
事实上，若 $L(s,\tau)=1$，则 $L(s,\sigma_i)=1$，由(2.7.1)有 $L(s,\xi)=1$，再由(2.7.4)得 $L(s,\pi)=1$。若 $L(s,\tau)\neq1$，则 $j=1,r=n,\tau=\alpha^t$，断言恰为 [R.G.--H.J.] 命题(7.11)。

\paragraph*{注 3.1.4。} 若 $\pi$ 平方可积，则 $L(s,\pi)$ 的极点位于半平面 $\operatorname{Re}s\leq0$。这由(3.1.3)或 [R.G.--H.J.] 命题(1.3)得到。

\origsubsection{3.2}{缓增表示（Tempered representations）}
仍令 $P$ 为类型 $(n_1,n_2,\ldots,n_r)$ 的抛物子群，现令 $\sigma_i$ 为 $G_{n_i}$ 的平方可积表示。众所周知，诱导表示
\begin{equation}
 \pi=I(G,P;\sigma_1,\sigma_2,\ldots,\sigma_r)\tag{3.2.1}
\end{equation}
不可约（例如参见 [H.J.2]）。此类不可约表示恰好就是缓增表示。注意，若 $\pi$ 等价于另一个表示
\begin{equation}
 \pi'=I(G,P';\sigma'_1,\sigma'_2,\ldots,\sigma'_{r'})\tag{3.2.2}
\end{equation}
且各 $\sigma'_i$ 平方可积，则 $P',P$ 相联系（associate）。更准确地，$r=r'$，$P=MU_P$，$P'=M'U_{P'}$，且存在 $G_n$ 的内自同构把 $M$ 映到 $M'$，把 $\sigma=\mathop{\times}_i\sigma_i$ 映到 $\sigma'=\mathop{\times}_i\sigma'_i$；换言之，从 $(M,(\sigma_i))$ 到 $(M',(\sigma'_i))$ 是置换 $M,M'$ 的各对角块。反过来，若 $P',\sigma'_i$ 以此方式与 $P,\sigma_i$ 联系，则 $\pi'$ 等价于 $\pi$。由(2.3)，
\begin{equation}
\begin{gathered}
 L(s,\pi)=\prod_iL(s,\sigma_i),\qquad L(s,\widetilde\pi)=\prod_iL(s,\widetilde\sigma_i),\\
 \varepsilon(s,\pi,\psi)=\prod_i\varepsilon(s,\sigma_i,\psi).
\end{gathered}\tag{3.2.3}
\end{equation}

\paragraph*{注 3.2.4（Remark 3.2.4）。}
由(3.2.3)及(3.1.4)，若 $\pi$ 缓增，则 $L(s,\pi)$ 的极点位于半平面 $\operatorname{Re}s\leq0$。

若表示 $\pi$ 形如 $\pi=\pi_0\otimes\alpha^t$，其中 $\pi_0$ 缓增，$t$ 为实数，则称它为本质缓增的（essentially tempered）。于是
\begin{equation}
 L(s,\pi)=L(s+t,\pi_0).\tag{3.2.5}
\end{equation}

% Jacquet，第3节后半；原书第71--75页。由23-jacquet-l-nonarch-factor.tex接入。
\origsubsection{3.3}{Langlands 构造（Langlands construction）}
下面回顾 Silberger 和 Wallach 将 Langlands 的结果推广到 $p$ 进情形的工作。

设 $Q=M_QU_Q$ 为型 $(p_1,p_2,\ldots,p_r)$ 的抛物子群（parabolic subgroup），并对每个 $1\leq i\leq r$，取 $G_{m_i}$ 的一个不可约本质缓增表示（irreducible essentially tempered representation）$\tau_i$。置 $\tau=\mathop{\times}_i\tau_i$。于是 $\tau_i=\tau_{i,0}\otimes\alpha^{t_i}$，其中 $\tau_{i,0}$ 缓增，$t_i$ 为实数。假定
\begin{equation}
t_1>t_2>\cdots>t_r.\tag{3.3.1}
\end{equation}
则诱导表示（induced representation）
\begin{equation}
\eta=I(G,Q;\tau_1,\tau_2,\ldots,\tau_r)\tag{3.3.2}
\end{equation}
有一个最大的真子表示 $\eta'$（它可以是 $\{0\}$）；不可约表示 $\pi=\eta/\eta'$ 记作
\begin{equation}
\pi=J(G,P;\tau_1,\tau_2,\ldots,\tau_r).\tag{3.3.3}
\end{equation}
每个不可约表示 $\pi$ 都有 (3.3.3) 的形式，其中（标准的）$P$ 及各 $\tau_i$ 由 $\pi$ 唯一确定。\footnote{译注：扫描页在本段先以 $p_i$ 表示分块大小，随后写作 $G_{m_i}$；式 (3.3.2) 用 $Q$，而式 (3.3.3) 及其后一句用 $P$。这里保留可见记号。}

$\eta'$ 的空间 $V'$ 可以显式描述。设 $W$ 为 $\tau$ 的空间，$\widetilde W$ 为逆反表示（contragredient representation）$\widetilde\tau$ 的空间。则 $V'$ 由 $\xi$ 的空间 $V$ 中满足下式的 $F$ 组成：
\begin{equation}
\int\langle F(\bar u g),\widetilde v\rangle\,d\bar u=0,
\qquad \widetilde v\in\widetilde W,\quad g\in G_n.\tag{3.3.4}
\end{equation}
此积分绝对收敛，积分域为 $\bar U=\bar U_Q=U_{\bar Q}$，其中 $\bar Q={}^tQ$ 是与 $Q$ 相反的抛物子群（opposite parabolic subgroup）。\footnote{译注：此处表示空间所属的表示在扫描页写作 $\xi$，前文相应诱导表示记为 $\eta$。}容易推出，$\pi$ 的矩阵系数（coefficient）可由与 (2.2.4) 类似的积分得到。具体地，设 $H:G_n\times G_n\longrightarrow\C$ 满足下列性质：
\begin{equation}
\begin{gathered}
H(u_1mg_1,\bar u_2mg_2)=H(g_1,g_2),\\
u_1\in U_Q,\quad\bar u_2\in\bar U_Q,\quad m\in M_Q;
\end{gathered}\tag{3.3.5}
\end{equation}
\begin{equation}
\begin{gathered}
\text{对任意 }g_1,g_2,\quad m\longmapsto H(mg_1,g_2)\\
\text{是 }\tau\otimes\delta_Q^{1/2}\text{ 的矩阵系数；}
\end{gathered}\tag{3.3.6}
\end{equation}
\begin{equation}
H\text{ 关于右作用是 }K_n\times K_n\text{ 有限的。}\tag{3.3.7}
\end{equation}
则由收敛积分
\begin{equation}
f(g)=\int_{M\backslash G}H(hg,h)\,dh
=\int_{\bar U\times K}H(\bar u kg,k)\,dk\,d\bar u\tag{3.3.8}
\end{equation}
定义的函数 $f$ 是 $\pi$ 的矩阵系数，并且 $\pi$ 的所有矩阵系数都可通过适当的 $H$ 以此方式得到（关于收敛问题，见 (3.6.6) 和 (3.6.7)）。

可以用 $\bar Q$ 代替 $Q$。此时，只要将条件 (3.3.1) 换成不等号方向相反的类似条件，就能定义与 (3.3.3) 相应的表示。例如，商表示
\begin{equation}
\pi'=J(G,\bar Q;\widetilde\tau_1,\widetilde\tau_2,\ldots,\widetilde\tau_r)\tag{3.3.9}
\end{equation}
是有定义的，它是诱导表示
\begin{equation}
\eta'=I(G,\bar Q;\widetilde\tau_1,\widetilde\tau_2,\ldots,\widetilde\tau_r)\tag{3.3.10}
\end{equation}
的商。其矩阵系数由积分 (3.3.8) 给出，其中 $H$ 满足 (3.1.5)--(3.1.7)，并将 $(Q,\tau)$ 换成 $(\bar Q,\widetilde\tau)$。\footnote{译注：扫描页此处交叉引用可见为 (3.1.5)--(3.1.7)；上文对应的三个条件编号为 (3.3.5)--(3.3.7)。}

特别地，设 $f$ 为由 (3.3.8) 定义的 $\pi$ 的矩阵系数，其中 $H$ 对 $(Q,\tau)$ 满足 (3.3.5)--(3.3.7)。则
\begin{equation}
f^\vee(g)=\int_{M\backslash G}\widetilde H(hg,h)\,dh,\tag{3.3.11}
\end{equation}
其中 $\widetilde H(g_1,g_2)=H(g_2,g_1)$。而 $\widetilde H$ 确实对 $(\bar Q,\widetilde\tau)$ 满足 (3.3.5)--(3.3.7)。所以 $f^\vee$ 是 $\pi'$ 的矩阵系数，且
\begin{equation}
\pi'=\widetilde\pi.\tag{3.3.12}
\end{equation}
当然，$\bar Q$ 也与型 $(n_r,n_{r-1},\ldots,n_1)$ 的标准抛物子群 $Q'$ 共轭，因而
\begin{equation}
\widetilde\pi=\pi'=J(G,Q';\widetilde\tau_r,\widetilde\tau_{r-1},\ldots,\widetilde\tau_1).\tag{3.3.13}
\end{equation}

\begin{originaltheorem}{定理 3.4（Theorem 3.4）}
沿用 (3.3) 的记号。则
\[
L(s,\pi)=\prod_iL(s,\tau_i),\qquad
L(s,\widetilde\pi)=\prod_iL(s,\widetilde\tau_i),
\]
\[
\varepsilon(s,\pi,\psi)=\prod_i\varepsilon(s,\tau_i,\psi).
\]
\end{originaltheorem}
\begin{proof}
只需证明关于 $L(s,\pi)$ 的断言。事实上，交换 $\pi$ 与 $\widetilde\pi$ 即可得到关于 $L(s,\widetilde\pi)$ 的断言；最后一个断言则由 (2.7) 和关系
\[
\gamma(s,\eta,\psi)=\prod_i\gamma(s,\tau_i,\psi)
\]
（参见 (2.3)）推出。证明分为本小节余下的论证和 (3.5) 中的构造两部分。由 (2.3) 和 (2.7)，存在多项式 $P$ 和 $\widetilde P$，使
\[
L(s,\pi)=P(q^{-s})\prod_iL(s,\tau_i),\qquad
L(s,\widetilde\pi)=\widetilde P(q^{-s})\prod_iL(s,\widetilde\tau_i).
\]
由 (3.2.4) 与 (3.3.1) 可知，$L(s,\tau_1)^{-1}$ 与 $\prod_iL(1-s,\widetilde\tau_i)^{-1}$ 互素。因此，$P(q^{-s})$ 与 $L(s,\tau_1)^{-1}$ 互素（(2.7.5)）。所以，若 $P\ne1$，则
\begin{equation}
\prod_{2\leq i\leq r}L(s,\tau_i)\tag{3.4.1}
\end{equation}
不属于 $I(\pi)$。于是只需证明 (3.4.1) 实际上属于 $I(\pi)$。

\paragraph*{3.5.}
恒等式
\begin{equation}
L(s,\pi)=\prod_{1\leq i\leq r}L(s,\tau_i)\tag{3.5.1}
\end{equation}
在 $r=1$ 时显然成立。假定它对 $r-1$ 成立，来证明它对 $r$ 成立。如刚才所见，只需证明 (3.4.1) 属于 $I(\pi)$。

置 $n_1=p_1$，$n_2=p_2+\cdots+p_r$。设 $Q'$ 为 $G_{n_2}$ 中型 $(p_2,\ldots,p_r)$ 的抛物子群。于是表示
\begin{equation}
\sigma_2=J(G_{n_2},Q';\tau_2,\ldots,\tau_r)\tag{3.5.2}
\end{equation}
是有定义的。再置 $\sigma_1=\tau_1$、$\sigma=\sigma_1\times\sigma_2$，设 $P$ 为 $G_n$ 中型 $(n_1,n_2)$ 的抛物子群。则 $\pi$ 实际上是
\begin{equation}
\xi=I(G,P;\sigma_1,\sigma_2)\tag{3.5.3}
\end{equation}
的一个商表示。更精确地，容易看出，$\pi$ 的矩阵系数由积分 (3.3.8) 给出，其中 $H$ 为对 $(P,\sigma)$ 满足 (3.3.5)--(3.3.7) 的任意函数，即用 $(P,\sigma)$ 代替 $(Q,\tau)$。

若 $f$ 是以这种方式定义的 $\pi$ 的矩阵系数，则
\begin{equation}
\begin{aligned}
&Z\bigl(\Phi,s+\tfrac12(n-1),f\bigr)\\
&\quad=\iint dk\,d\bar u\int_{G_n}|\det g|^{s+(n-1)/2}
\Phi(k^{-1}\bar u^{-1}g)H(g,k)\,d^\times g,
\end{aligned}\tag{3.5.4}
\end{equation}
其中 $\bar u$ 的积分域为 $\bar U=U_{\bar P}$。这也等于
\begin{equation}
\begin{aligned}
&\iint dk\,dk'\iint
|\det m_1|^{s+(n_1-1)/2}|\det m_2|^{s+(n_2-1)/2}\\
&\quad\cdot\delta_{\bar P}\!\begin{pmatrix}m_1&0\\0&m_2\end{pmatrix}^{-1/2}
H\!\left[\begin{pmatrix}m_1&0\\0&m_2\end{pmatrix}k',k\right]
\,d^\times m_1\,d^\times m_2\\
&\quad\cdot\iint\Phi\!\left[k^{-1}\begin{pmatrix}m_1&y\\xm_1&m_2+xy\end{pmatrix}k'\right]\,dx\,dy.
\end{aligned}\tag{3.5.5}
\end{equation}
这里 $m_i$ 遍历 $G_{n_i}$。我们将看到，给定 $\Phi_2\in\mathcal S(n_2\times n_2,F)$ 和 $\sigma_2$ 的矩阵系数 $f_2$，存在 $H$ 和 $\Phi$，使 $Z(\Phi_2,s+\tfrac12(n_2-1),f_2)$ 等于 (3.5.5)。由于根据归纳假设，$L(s,\sigma_2)$ 等于 (3.4.1)，由此便有 (3.4.1) 属于 $I(\pi)$，证明也就完成了。

存在 $\pi_1$ 的一个矩阵系数 $f_1$ 和一个 $\Phi_1\in\mathcal S(n_1\times n_1,F)$，后者的支集包含在 $e$ 的某个紧邻域内，使得
\begin{equation}
Z\bigl(\Phi_1,s+\tfrac12(n_1-1),f_1\bigr)=1.\tag{3.5.6}
\end{equation}
还存在 $\Phi_{12}\in\mathcal S(n_1\times n_2,F)$ 和 $\Phi_{21}\in\mathcal S(n_2\times n_1,F)$，它们的支集包含在 $0$ 的某个邻域内，使得由
\begin{equation}
\Phi\!\begin{pmatrix}m_1&y\\x&m_2\end{pmatrix}
=\Phi_1(m_1)\Phi_{12}(y)\Phi_{21}(x)\Phi_2(m_2)\tag{3.5.7}
\end{equation}
定义的 $\Phi\in\mathcal S(n\times n,F)$ 满足
\begin{equation}
\iint\Phi\!\begin{pmatrix}m_1&y\\xm_1&m_2+xy\end{pmatrix}\,dx\,dy
=\Phi_1(m_1)\Phi_2(m_2).\tag{3.5.8}
\end{equation}
还存在 $K$ 上两个 $K$ 有限函数（$K$-finite functions）$\xi$ 和 $\xi'$，使得
\begin{equation}
\iint\Phi(k^{-1}zk')\xi(k)\xi'(k')\,dk\,dk'=\Phi(z).\tag{3.5.9}
\end{equation}
于是
\begin{equation}
\begin{aligned}
&Z\bigl(\Phi_2,s+\tfrac12(n-1),f_2\bigr)\\
&\quad=\iint|\det m_1|^{s+(n_1-1)/2}|\det m_2|^{s+(n_2-1)/2}
 f_1(m_1)f_2(m_2)\,d^\times m_1\,d^\times m_2\\
&\qquad\cdot\iint\xi(k)\xi'(k')\,dk\,dk'
\iint\Phi\!\left[k^{-1}\begin{pmatrix}m_1&y\\xm_1&m_2+xy\end{pmatrix}k'\right]\,dx\,dy.
\end{aligned}\tag{3.5.10}
\end{equation}


最后沿用 (2.5) 的平均化方法完成证明。\footnote{译注：扫描页在引入 $f_1$ 时写 $\pi_1$，本段上文相应表示为 $\sigma_1$；式 (3.5.10) 左端写 $s+\tfrac12(n-1)$，而前文同一 $\Phi_2,f_2$ 的积分写作 $s+\tfrac12(n_2-1)$。式 (3.5.5) 的模特征下标可见为 $\bar P$，而式 (3.5.11) 为 $P$。均按可见原式保留。}具体地，设 $dh$ 和 $dh'$ 分别为 $K\cap\bar P$ 和 $K\cap P$ 上的归一化 Haar 测度（normalized Haar measure）。对 $h\in K\cap\bar P$、$h'\in K\cap P$，写
\[
h=\begin{pmatrix}h_1&0\\ *&h_2\end{pmatrix},\qquad
h'=\begin{pmatrix}h'_1&*\\0&h'_2\end{pmatrix}.
\]
则 (3.5.10) 还等于
\begin{equation}
\begin{aligned}
&\iint dk\,dk'\iint|\det m_1|^{s+(n_1-1)/2}|\det m_2|^{s+(n_2-1)/2}\\
&\quad\cdot\delta_P\!\begin{pmatrix}m_1&0\\0&m_2\end{pmatrix}^{-1/2}
H_1\!\left[\begin{pmatrix}m_1&0\\0&m_2\end{pmatrix},k,k'\right]
\,d^\times m_1\,d^\times m_2\\
&\quad\cdot\iint\Phi\!\left[k^{-1}\begin{pmatrix}m_1&y\\xm_1&m_2+xy\end{pmatrix}k'\right]\,dx\,dy,
\end{aligned}\tag{3.5.11}
\end{equation}
其中
\begin{equation}
\begin{aligned}
&H_1\!\left[\begin{pmatrix}m_1&0\\0&m_2\end{pmatrix},k,k'\right]\\
&\quad=\delta_P\!\begin{pmatrix}m_1&0\\0&m_2\end{pmatrix}^{1/2}
\iint\xi(hk)\xi'(h'k')f_1(h_1m_1h_1'^{-1})f_2(h_2m_2h_2'^{-1})\,dh\,dh'.
\end{aligned}\tag{3.5.12}
\end{equation}
存在对 $(P,\sigma)$ 满足 (3.3.5)--(3.3.7) 的 $H$，使
\begin{equation}
H\!\left[\begin{pmatrix}m_1&0\\0&m_2\end{pmatrix}k',k\right]
=H_1\!\left[\begin{pmatrix}m_1&0\\0&m_2\end{pmatrix},k,k'\right].\tag{3.5.13}
\end{equation}
比较 (3.5.11) 与 (3.5.5)，便得到所要的断言。
\end{proof}

\origsubsection{3.6}{非分歧表示（Unramified representations）}
设 $\pi$ 非分歧，即含有非零的 $K$-不变向量（也即包含 $K$ 的平凡表示，unit representation）。用 $B_n$ 表示型 $(1,1,1,\ldots,1)$ 的抛物子群，则存在 $u_i\in\C$，使 $\pi$ 为下列诱导表示的唯一非分歧成分：
\begin{equation}
\xi=I(G,B_n;\mu_1,\mu_2,\ldots,\mu_n),\qquad \mu_i=\alpha^{u_i}.\tag{3.6.1}
\end{equation}
甚至可以假定 $u_1\geq u_2\geq\cdots\geq u_n$。改变记号，可将由各 $\mu_i$ 组成的 $n$ 元组写成
\[
(\nu_1^1,\nu_2^1,\ldots,\nu_{n_1}^1,
 \nu_1^2,\nu_2^2,\ldots,\nu_{n_2}^2,\ldots,
 \nu_1^r,\nu_2^r,\ldots,\nu_{n_r}^r),
\]
其中 $n=n_1+n_2+\cdots+n_r$，$\nu_i^j=\alpha^{t_j}\lambda_i^j$，$\lambda_i^j$ 是特征（character），$t_j$ 为实数，且 $t_1>t_2>\cdots>t_r$。于是
\[
\tau_j=I(G_{n_j},B_{n_j};\nu_1^j,\nu_2^j,\ldots,\nu_{n_j}^j)
\]
是 $G_{n_j}$ 的本质缓增非分歧表示；事实上，若 $Q$ 为型 $(n_1,n_2,\ldots,n_r)$ 的抛物子群，则
\begin{equation}
\pi=J(G,Q;\tau_1,\tau_2,\ldots,\tau_r).\tag{3.6.2}
\end{equation}
这由一个显式计算得出：若 $H$ 对 $(Q,\mathop{\times}_j\tau_j)$ 满足条件 (3.3.5)--(3.3.7)，且 $H(k,k')=1$，则
\begin{equation}
\int H(\bar u,e)\,d\bar u\ne0.\tag{3.6.3}
\end{equation}
所以，$\eta=I(G,Q;\tau_1,\tau_2,\ldots,\tau_r)$ 的“$J$ 成分”是非分歧的；由于 $\xi\simeq\eta$，便得到 (3.6.2)。\footnote{译注：扫描页先取 $u_i\in\C$，随后直接写 $u_1\geq\cdots\geq u_n$，未在该行另写实部；此处保留其可见形式。}

因此
\begin{equation}
\begin{gathered}
L(s,\pi)=\prod_iL(s,\mu_i),\qquad
L(s,\widetilde\pi)=\prod_iL(s,\mu_i^{-1}),\\
\varepsilon(s,\pi,\psi)=\prod_i\varepsilon(s,\mu_i,\psi)
\quad\text{（若 $\psi$ 的指数为零，则等于 $1$）。}
\end{gathered}\tag{3.6.4}
\end{equation}
更精确的结果见 [R.G.-H.J., \S6]。

\begin{originaltheorem}{备注 3.6.5（Remark 3.6.5）}
沿用 (2.3) 的记号。设各 $\sigma_i$ 不可约且非分歧。则 $\xi$ 有唯一的非分歧成分 $\pi$。由 (3.6.4) 容易推出
\[
L(s,\pi)=\prod_iL(s,\sigma_i),
\]
而 $L(s,\widetilde\pi)$ 和 $\varepsilon$ 也有类似的关系。
\end{originaltheorem}

\begin{originaltheorem}{备注 3.6.6（Remark 3.6.6）}
设 $\pi$ 是 $G_n$ 的缓增表示。令 $\rho_n=I(G_n,B_n;1,1,\ldots,1)$，以 $\Xi_n$ 表示附属于 $\rho$ 的球函数（spherical function）。则 $\pi$ 的每个矩阵系数的绝对值都以 $\Xi_n$ 的某个常数倍为上界。再沿用 (3.3) 的记号，则 $\tau_i$ 的矩阵系数的绝对值都以 $\Xi_{n_i}\otimes\alpha^{t_i}$ 的某个常数倍为上界。此外，若 $H$ 对 $(Q,\tau)$ 满足 (3.3.5)--(3.3.7)，则 $|H|\leq cH_0$，其中 $H_0$ 满足 (3.3.5)，且
\[
H_0(mk,k')=\delta_P^{1/2}(m)\prod_i\Xi_{n_i}(m_i).
\]
所以，就收敛问题而言，可以将 $\tau_i$ 换成 $\rho_{n_i}\otimes\alpha^{t_i}$，并将 $H$ 换成 $H_0>0$。结合下一个引理，就解决了所有收敛问题。\footnote{译注：本显示式在扫描页的下标可见为 $P$，且乘积内未另写前文的 $\alpha^{t_i}$ 因子；这里保留原式。}
\end{originaltheorem}

\begin{originaltheorem}{引理 3.6.7（Lemma 3.6.7）}
若 $\Omega$ 是 $\mathcal S(n\times n,F)$ 的一个有界集（bounded set），则存在 $\mathcal S(n\times n,F)$ 中的 $\Phi_0\geq0$，使得对 $k_i\in K$ 有 $\Phi_0(k_1xk_2)=\Phi_0(x)$，且对 $\Phi\in\Omega$ 有 $|\Phi|\leq\Phi_0$。
\end{originaltheorem}
\begin{proof}
可以假定，$\Omega$ 还包含所有由 $\Phi\in\Omega$ 得到的函数 $x\mapsto\Phi(k_1xk_2)$。于是存在 $\Phi_1\geq0$，使得对 $\Omega$ 中所有的 $\Phi$ 都有 $|\Phi|\leq\Phi_1$ [A.W. 1, Lemma 5]。只需取
\[
\Phi_0(x)=\iint\Phi_1(k_1xk_2)\,dk_1\,dk_2,
\]
其中 $dk_i$ 是 $K$ 上的归一化 Haar 测度。
\end{proof}

\paragraph*{3.7.}
按照“Langlands 哲学（Langlands' philosophy）”，Weil--Deligne 群（Weil--Deligne group）$W'_F$ 的 $n$ 维半单表示（semisimple representation）与 $G_n(F)$ 的不可约容许表示（irreducible admissible representation）之间，应当存在一个“自然双射（natural bijection）”$\sigma\mapsto\pi(\sigma)$。此外，对 $\pi=\pi(\sigma)$，还应有
\[
L(s,\pi)=L(s,\sigma),\qquad
\varepsilon(s,\pi,\psi)=\varepsilon(s,\sigma,\psi),
\]
\[
\pi(\sigma)^\sim=\pi(\widetilde\sigma),\qquad
\pi(\sigma\otimes\chi)=\pi(\sigma)\otimes\chi.
\]
显然，若能对 Weil 群（Weil group）的不可约表示定义映射 $\sigma\mapsto\pi(\sigma)$，则这一猜想就能得到证明（参见下文第5节）。


% 原书75页第4节起至86页末；主文件由E组持有并input本片段。
\origsection{4}{局部 Archimedes 理论（Local archimedean theory）}
第 4、5 节均设基域 $F$ 为 Archimedes 局部域（archimedean local field）。

\paragraph*{4.1.}
当 $F=\R$ 时令 $K_n=O(n,\R)$，当 $F=\C$ 时令 $K_n=U(n)$。用 $\mathfrak g_n$ 表示实 Lie 群 $G_n(F)$ 的 Lie 代数（Lie algebra）。这里只考虑 $(\mathfrak g_n,K_n)$ 的容许表示（admissible representation），定义见 [N.W.]；为叙述方便，仍常称其为 $G_n(F)$ 的表示。以下沿用群表示与 $(\mathfrak g_n,K_n)$-模之间的通常联系。

设 $\pi$ 为 $(\mathfrak g_n,K_n)$ 的容许表示。即使 $\pi$ 不是群的表示，也可以定义其逆反表示（contragredient representation）$\widetilde\pi$、矩阵系数（coefficient）及中心拟特征（central quasi-character），参见[H.J.-R.L., §§5、6]。与前面一样，若 $f$ 为 $\pi$ 的系数，则 $f^\vee$（1.1.2）为 $\widetilde\pi$ 的系数。

与非 Archimedes 情形一样，令 $\mathcal S(p\times q,F)$ 为 $M(p\times q,F)$ 上的 Schwartz 函数（Schwartz function）空间。引入子空间 $\mathcal S_0=\mathcal S_0(p\times q,F)$，由下列形式的函数 $\Phi$ 构成，其中 $P$ 为多项式：
\[
\begin{aligned}
 \Phi(x)&=P(x_{ij})\exp(-\pi\textstyle\sum x_{ij}^2),&&F=\R,\\
 \Phi(z)&=P(z_{ij},\overline z_{ij})\exp(-2\pi\textstyle\sum z_{ij}\overline z_{ij}),&&F=\C.
\end{aligned}
\]
给定 $\pi$，仍可考察 (1.1.3) 中的积分，其中 $\Phi\in\mathcal S(p\times q,F)$，$f$ 为 $\pi$ 的系数。\footnote{\textbf{译者注（本次校订）：}此处须取 $p=q=n$，即 $\Phi\in\mathcal S(n\times n,F)$，因为 (1.1.3) 的积分变量是 $n$ 阶可逆矩阵。}

\begin{originaltheorem}{命题 4.2（Proposition 4.2）}
设 $\pi$ 为 $(\mathfrak g_n,K_n)$ 的不可约容许表示。

(1) 存在 $s_0$，使积分(1.1.3)在 $\Re s>s_0$ 时绝对收敛。

(2) 对 $\Phi\in\mathcal S_0(n\times n,F)$，积分是全复平面上 $s$ 的亚纯函数（meromorphic function）。更确切地说，可写为 $s$ 的多项式乘以某个固定亚纯函数。

(3) 设 $\psi\ne1$ 为 $F$ 的加法特征（additive character）。存在亚纯函数 $\gamma(s,\pi,\psi)$，使对 $\pi$ 的全部系数 $f$ 及全部 $\Phi\in\mathcal S_0(n\times n,F)$，有
\[
 Z(\Phi^\wedge,1-s+\tfrac12(n-1),f^\vee)
 =\gamma(s,\pi,\psi)Z(\Phi,s,f).\tag{4.2.3}
\]
\end{originaltheorem}
$\Phi^\wedge$ 仍由 (1.2.4) 定义。\footnote{\textbf{译者注（本次校订）：}与 (1.2.3) 一样，(4.2.3) 右端的积分参数应为 $s+\tfrac12(n-1)$；这样才能与本节后续的 $L$ 因子定义及 (1.3.7) 一致。}若 $\Phi\in\mathcal S_0$，而 $\psi$ 由下式定义，
\[
\begin{aligned}
 \psi(x)&=\exp(\pm2i\pi x),&&F=\R,\\
 \psi(z)&=\exp[\pm2i\pi(z+\overline z)],&&F=\C,
\end{aligned}\tag{4.2.4}
\]
则 $\Phi^\wedge$ 仍属于 $\mathcal S_0$；由(4.2.2)应用于 $\widetilde\pi$，(4.2.3)左边有定义。若 $\psi$ 不由(4.2.4)给出，则需相应调整正规化，参见(1.3.11)、(1.3.12)。以下假设 $\psi$ 由(4.2.4)给出。

\paragraph*{例 4.2.5（Example 4.2.5）.}
若 $n=1$，$\pi$ 就是 $F^\times$ 的拟特征，命题4.2已在[J.T.]或[A.W.2]中证明。

\paragraph*{4.3.}
命题 4.2 的证明见 4.4。这里先推导几个直接推论。由引理3.6.7，(1.1.3)在 $\Re s>s_0$ 时对有界集中的 $\Phi$ 一致收敛；因此
\[
 \Phi\longmapsto Z(\Phi,s,f)\tag{4.3.1}
\]
是在 $\Re s>s_0$ 上全纯依赖于 $s$ 的分布（distribution）。由于 $\mathcal S_0(n\times n,F)$ 在 $\mathcal S(n\times n,F)$ 中稠密，若 $f\ne0$，则至少存在 $\Phi\in\mathcal S_0$，使 $Z(\Phi,s,f)\not\equiv0$。

假设(4.2.2)已证。取亚纯函数 $L(s)$，使所有 $\Phi\in\mathcal S_0$ 的比值 $Z(\Phi,s+\tfrac12(n-1),f)/L(s)$ 均为多项式；用 $I$ 表示这些比值在 $\C[s]$ 中张成的复向量子空间。由前面的讨论，$I\ne\{0\}$。若 $\Phi\in\mathcal S_0$，则下式定义的 $\Phi'$ 也属于 $\mathcal S_0$：
\[
 \Phi'(x)=\left.\frac d{dt}\Phi(xe^{-t})\right|_{t=0}.\tag{4.3.2}
\]
设 $\omega$ 为 $\pi$ 的中心拟特征。对恒等式
\begin{equation}
\begin{split}
 Z(\Phi,s+\tfrac12(n-1),f)
 &=\int\Phi(xe^{-t})f(x)|\det x|^{s+(n-1)/2}\,d^\times x\\
 &\quad\cdot\omega(e^{-t})\exp[-ts-\tfrac t2(n-1)]
\end{split}\tag{4.3.3}
\end{equation}
关于 $t$ 求导，再令 $t=0$，得到形如\footnote{\textbf{译者注（本次校订）：}(4.3.3) 中的指数因子漏了行列式的维数。令 $d=[F:\R]$，则 $|\det(e^{-t}1_n)|_F=e^{-dnt}$，故该因子应为 $\exp[-dnt(s+\tfrac12(n-1))]$，同时保留 $\omega(e^{-t})$。修正只改变 (4.3.4) 中的常数 $a,b$，其中 $a$ 仍非零。}
\[
 (as+b)Z(\Phi,s+\tfrac12(n-1),f)
 +Z(\Phi',s+\tfrac12(n-1),f)=0,\qquad a\ne0.\tag{4.3.4}
\]
的关系。这表明 $I$ 为理想（ideal）。取 $I$ 的生成元 $P_0$，令 $L(s,\pi)=P_0(s)L(s)$。这些比值仍为多项式，但这次它们作为向量空间张成整个 $\C[s]$。精确到非零常数乘因子，这些性质刻画了 $L(s,\pi)$。

现假设命题 4.2 对 $(\pi,\widetilde\pi)$ 成立。则 $\gamma$ 唯一确定；按 (1.3.6) 定义 $\varepsilon$ 因子后，函数方程 (4.2.3) 就可写成 (1.3.7) 的形式。由 $L$ 因子的定义，这蕴含 $\varepsilon(s,\pi,\psi)$ 及其逆均属于 $\C[s]$。所以当 $\psi$ 由(4.2.4)给出时，它只是常数；否则为指数因子（1.3.9）。

\paragraph*{注 4.3.5.}
目前 $L$、$\varepsilon$ 因子都只确定到一个非零常数倍。由于 $\gamma$ 已由积分的函数方程唯一确定，这些常数之间也有相应的约束。$n=1$ 时可取[J.T.1、2]中给出的因子。
\paragraph*{注 4.3.6.}
关系(1.3.9)--(1.3.12)适用于 Archimedes 情形。

\paragraph*{4.4.}
记号如2.1，而基域改为 $\R$ 或 $\C$；$\sigma_i$ 此时自然为 $(\mathfrak g_{n_i},K_{n_i})$ 的容许表示。仍可定义诱导表示（induced representation）$\xi$（2.1.4），见[N.W.]；其系数由2.2的规则给出，区别仅在于(2.2.2)中 $g_i$ 必须属于 $K_n$，(2.2.3)中 $H$ 须在 $G_n\times G_n$ 上为 $C^\infty$。2.3之前的说明仍适用，与2.3一样，有：

\begin{originaltheorem}{命题 4.4.1（Proposition 4.4.1）}
沿用2.1记号，假设每个 $\sigma_i$ 有中心拟特征，且4.2对每对 $(\sigma_i,\widetilde\sigma_i)$ 成立。则它也对 $(\xi,\widetilde\xi)$ 成立。此外，$\gamma(s,\xi,\psi)=\prod_i\gamma(s,\sigma_i,\psi)$，且可取
\[
 L(s,\xi)=\prod_iL(s,\sigma_i),\qquad
 L(s,\widetilde\xi)=\prod_iL(s,\widetilde\sigma_i),
 \qquad\varepsilon(s,\xi,\psi)=\prod_i\varepsilon(s,\sigma_i,\psi).
\]
\end{originaltheorem}
证明沿用 2.3 的方法，只需注意：若 $\Phi\in\mathcal S_0(n\times n,F)$，则 $x\mapsto\Phi(k_1xk_2)$（$k_i\in K_n$）张成 $\mathcal S_0$ 的有限维子空间；对每个 $k,k'\in K_n$，函数(2.4.2)属于 $\bigotimes_i\mathcal S_0(n_i\times n_i,F)$。

设 $\pi$ 为 $\xi$ 的一个不可约组成因子（irreducible component）。$\pi$ 的每个系数都是 $\xi$ 的系数，且 $\widetilde\pi$ 为 $\widetilde\xi$ 的组成因子，所以4.2对 $(\pi,\widetilde\pi)$ 成立。更精确地，存在多项式 $R,\widetilde R$，使
\[
 L(s,\pi)=R(s)\prod_iL(s,\sigma_i),\qquad
 L(s,\widetilde\pi)=\widetilde R(s)\prod_iL(s,\widetilde\sigma_i),\tag{4.4.2}
\]
而
\[
 \gamma(s,\pi,\psi)=\prod_i\gamma(s,\sigma_i,\psi).\tag{4.4.3}
\]
因为 $\varepsilon(s,\pi,\psi)$、$\varepsilon(s,\sigma_i,\psi)$ 是常数，$R,\widetilde R$ 具有形式
\[
 R(s)=\prod_i(s-s_i),\qquad
 \widetilde R(s)=c\prod_i(1-s-s_i).\tag{4.4.4}
\]
最后，$G_n$ 的任意不可约容许表示 $\pi$ 都是4.4.1中某个 $n_i=1$ 的诱导表示 $\xi$ 的组成因子。此时4.2对每个 $(\sigma_i,\widetilde\sigma_i)$ 成立，故对 $(\pi,\widetilde\pi)$ 成立，4.2证毕。

\begin{originaltheorem}{命题 4.5（Proposition 4.5）}
对所有 $s$，$L(s,\pi)\ne0$。
\end{originaltheorem}
\begin{proof}
设 $\pi$ 为 $G_n$ 的不可约容许表示，作为4.4.1中 $n_i=1$ 的某个诱导表示 $\xi$ 的组成因子。先证明对任意 $\Phi\in\mathcal S(n\times n,F)$，比值
\[
 Z(\Phi,s+\tfrac12(n-1),f)/L(s,\pi)\tag{4.5.1}
\]
可延拓为 $s$ 的整函数（entire function）。

取 $s_0$，使 $Z(\Phi,s+\tfrac12(n-1),f)$ 在 $\Re s>s_0$ 时绝对收敛，且 $Z(\Phi,1-s+\tfrac12(n-1),f^\vee)$ 在 $\Re s<1-s_0$ 时绝对收敛。选取 $s_0<s_1<s_2$，再取多项式 $Q$，使 $Q(s)L(s,\xi)$ 在带形域 $1-s_2\leq\Re s\leq s_2$ 中无极点。令 $\Phi_i\in\mathcal S_0$ 收敛到 $\Phi\in\mathcal S$。对于给定 $f$ 及多项式 $R$，以下极限在 $s_1\leq\Re s\leq s_2$ 上一致成立：
\[
 \lim_{i\to\infty}R(s)Z(\Phi_i,s+\tfrac12(n-1),f)
 =R(s)Z(\Phi,s+\tfrac12(n-1),f),\tag{4.5.2}
\]
\[
 \lim_{i,j\to\infty}R(s)Z(\Phi_i-\Phi_j,s+\tfrac12(n-1),f)=0.\tag{4.5.3}
\]
同样，$\Phi_i^\wedge-\Phi_j^\wedge\longrightarrow0$，故在 $1-s_2\leq\Re s\leq1-s_1$ 上一致有
\[
 \lim_{i,j\to\infty}R(s)Z(\Phi_i^\wedge-\Phi_j^\wedge,
 1-s+\tfrac12(n-1),f^\vee)=0.\tag{4.5.4}
\]
反复使用(4.3.4)可将这些断言归约到 $R=1$，此时显然成立。另一方面，
\begin{align*}
 &Q(s)Z(\Phi_i-\Phi_j,s+\tfrac12(n-1),f)\\
 &=\varepsilon(s,\xi,\psi)^{-1}\frac{L(s,\xi)}{L(1-s,\widetilde\xi)}
 Q(s)Z(\Phi_i^\wedge-\Phi_j^\wedge,1-s+\tfrac12(n-1),f^\vee).
\end{align*}
经典公式
\[
 \frac{\Gamma(x+iy)}{\Gamma(x'-iy)}\simeq |y|^{x-x'}
 \qquad(|y|\longrightarrow\infty)
\]
表明，$Q(s)L(s,\xi)/L(1-s,\widetilde\xi)$ 的绝对值在上述带形域中以某个多项式为上界。\footnote{\textbf{译者注（本次校订）：}上式是关于绝对值的渐近公式，应读作 $|\Gamma(x+iy)/\Gamma(x^{\prime}-iy)|\sim|y|^{x-x^{\prime}}$（固定实数 $x,x^{\prime}$）。复数比值本身还带有随 $y$ 变化的相位；以下估计只需要其绝对值。}因此在 $1-s_2\leq\Re s\leq1-s_1$ 上也一致有
\[
 \lim_{i,j\to\infty}Q(s)Z(\Phi_i-\Phi_j,s+\tfrac12(n-1),f)=0.\tag{4.5.5}
\]
固定 $f$，给定 $\epsilon>0$，选择 $N$，使 $i,j\geq N$ 且 $s_1\leq\Re s\leq s_2$ 或 $1-s_2\leq\Re s\leq1-s_1$ 时，
$|Q(s)Z(\Phi_i-\Phi_j,s+\tfrac12(n-1),f)|\leq\epsilon$。但
\[
 Q(s)Z(\Phi_i-\Phi_j,s+\tfrac12(n-1),f)=Q(s)R(s)L(s,\xi),
\]
其中 $R$ 为另一个多项式。$L(s,\xi)$ 在任意竖直带形域中迅速衰减，所以给定 $(i,j)$，存在 $M$，当 $|\Im s|\geq M$ 时，
\[
 |Q(s)Z(\Phi_i-\Phi_j,s+\tfrac12(n-1),f)|\leq\epsilon.\tag{4.5.6}
\]
由最大值原理（maximum principle），当 $i,j\geq N$ 且 $1-s_2\leq\Re s\leq s_2$ 时，(4.5.6)成立。因此 $Q(s)Z(\Phi_i,s+\tfrac12(n-1),f)$ 对该带形域上的一致收敛拓扑为 Cauchy 列。极限 $Z(s)$ 在 $1-s_2<\Re s<s_2$ 上全纯，且在 $s_1\leq\Re s<s_2$ 上等于 $Q(s)Z(\Phi,s+\tfrac12(n-1),f)$。$s_2$ 可以任意增大，由此得到 (4.5.1) 的全纯性。\footnote{\textbf{译者注（本次校订）：}这里还需补充除去 $Q(s)L(s,\pi)$ 零点的步骤。各 $Z(\Phi_i,s+\tfrac12(n-1),f)/L(s,\pi)$ 都是整函数。在任一给定点附近，选取边界上 $Q(s)L(s,\pi)$ 无零极点的小圆盘；上述一致收敛先给出这些整函数在圆周上的一致收敛，再由最大模原理得到圆盘内的一致收敛。因此极限商在中心点也全纯。}

现假设 $L(s_0,\pi)=0$，且 $f\ne0$。则对每个 $\Phi$，亚纯函数 $Z(\Phi,s+\tfrac12(n-1),f)$ 在 $s_0$ 为零。但若 $\Phi$ 紧支且支集包含于 $G_n$，该函数对所有 $s$ 由收敛积分(1.1.3)定义。因此，对 $s=s_0$ 及全部 $\Phi\in C_c^\infty(G_n)$，(1.1.3)为零，故 $f=0$，矛盾。所以 $L(s_0,\pi)\ne0$。
\end{proof}

\begin{originaltheorem}{推论 4.5.8（Corollary 4.5.8）}
设 $\xi$ 如4.4.1，$\pi$ 为其不可约组成因子，$R,\widetilde R$ 如(4.4.2)。若 $s_i$ 为 $R$（分别地，$\widetilde R$）的 $m$ 阶零点，则它是 $L(s,\xi)$（分别地，$L(s,\widetilde\xi)$）的至少 $m$ 阶极点。
\end{originaltheorem}

\origsection{5}{$L$ 因子的计算：Archimedes 情形（Computation of the $L$-factor: archimedean case）}
回忆，Weil 群 $W_F$ 的 $n$ 维半单表示（semisimple representation）的类，与 $G_n(F)$ 的不可约容许表示之间，有“自然双射”$\lambda\mapsto\pi(\lambda)$，见[R.L.]、[A.K.-G.Z.]、[N.W.]。特别地，$\pi(\lambda)^\sim=\pi(\widetilde\lambda)$。

\begin{originaltheorem}{定理 5.1（Theorem 5.1）}
对任意 $n$ 维 $\lambda$，令 $\pi=\pi(\lambda)$。则
\[
 L(s,\pi)=L(s,\lambda),\qquad
 L(s,\widetilde\pi)=L(s,\widetilde\lambda),\qquad
 \varepsilon(s,\pi,\psi)=\varepsilon(s,\lambda,\psi).
\]
\end{originaltheorem}
第 5 节将专门证明这一结果。由于 $L$ 因子尚未归一化，更准确的说法是，可以取 $L(s,\pi)=L(s,\lambda)$，且
\[
 \gamma(s,\pi,\psi)=\varepsilon(s,\lambda,\psi)
 L(1-s,\widetilde\lambda)/L(s,\lambda).
\]

\paragraph*{5.2. 平方可积表示（Square-integrable representations）.}
若 $\lambda$ 不可约，则 $\pi(\lambda)$ 本质平方可积（essentially square-integrable），反之亦然。更确切地，$n=1$ 时 $\lambda$ 是 $F^\times$ 的拟特征，$\lambda=\pi$，断言由 $\lambda$ 的 $L$、$\varepsilon$ 因子定义得出。若 $n\ne1$，则 $n=2$，断言已在[H.J.-R.L., 定理13.1、引理13.24与5.17]证明；事实上，这一情形的证明是对 4.4.1 中方法的进一步细化。

\paragraph*{5.3. 缓增表示（Tempered representations）.}
若 $\lambda$ 酉（unitary），则 $\pi$ 缓增，反之亦然。更精确地，$\lambda=\bigoplus\lambda_i$，各 $\lambda_i$ 不可约且维数为 $n_i$。令 $\sigma_i=\pi(\lambda_i)$。则
$\pi=I(G,P;\sigma_1,\ldots,\sigma_r)$，其中 $P$ 为型 $(n_1,\ldots,n_r)$ 的抛物子群（parabolic subgroup）。由4.4.1，
\[
 L(s,\pi)=\prod_iL(s,\sigma_i)=\prod_iL(s,\lambda_i)=L(s,\lambda).\tag{5.3.1}
\]
$L(s,\widetilde\pi)$ 与 $\varepsilon(s,\pi,\psi)$ 类似计算，断言随之成立。本质缓增（essentially tempered）情形由(1.3.11)、(1.3.12)及下式得出：
\[
 \pi(\lambda\otimes\chi)=\pi(\lambda)\otimes\chi.\tag{5.3.2}
\]
\paragraph*{注 5.3.3.}
由(5.3.1)，当 $\pi$ 缓增时，$L(s,\pi)$ 的极点位于 $\Re s\leq0$。

\paragraph*{5.4. 一般情形（General case）.}
一般地，可写 $\lambda=\bigoplus_i\mu_i$，$\mu_i=\mu_{i,0}\otimes\alpha^{t_i}$，其中 $\mu_{i,0}$ 酉、维数为 $p_i$，$t_i$ 为实数。令 $\tau_{i,0}=\pi(\mu_{i,0})$、$\tau_i=\pi(\mu_i)$；由(5.3.2)，$\tau_i=\tau_{i,0}\otimes\alpha^{t_i}$。可假设(3.3.1)成立。若 $Q$ 为型 $(p_1,\ldots,p_r)$ 的抛物子群，则(3.3.2)的诱导表示 $\eta$ 有唯一不可约商，记号如(3.3.3)；该商就是 $\pi=\pi(\lambda)$。由4.4.1与5.3已知
\[
 \gamma(s,\pi,\psi)=\prod_i\gamma(s,\tau_i,\psi)
 =\varepsilon(s,\lambda,\psi)L(1-s,\widetilde\lambda)/L(s,\lambda).
\]
所以只须计算 $L(s,\pi)$。交换 $\lambda,\widetilde\lambda$ 后可得到 $L(s,\widetilde\pi)$（因为 $\widetilde\pi=\pi(\widetilde\lambda)$），再从 $\gamma$ 得到 $\varepsilon$。只须证明 $L(s,\pi)=\prod_{1\leq i\leq r}L(s,\tau_i)$。$r=1$ 显然，可假设 $r>1$，且对 $r-1$ 已证。

由已经证明的结果，先可写成
\[
 L(s,\pi)=P(s)\prod_iL(s,\tau_i),\qquad
 L(s,\widetilde\pi)=\widetilde P(s)\prod_iL(s,\widetilde\tau_i),
\]
其中 $P,\widetilde P$ 为满足(4.4.4)关系的多项式。设 $s_0$ 为 $P$ 的 $u$ 阶零点。由4.5.8，它是 $\prod_iL(s,\tau_i)$ 的至少 $u$ 阶极点；又由(4.4.4)，它是 $\prod_iL(1-s,\widetilde\tau_i)$ 的至少 $u$ 阶极点。但 $L(s,\tau_1)$ 与 $\prod_iL(1-s,\widetilde\tau_i)$ 不可能有公共极点，见(5.3.2)、(5.3.3)、(3.3.1)。因此只须证明：$\prod_{i\geq2}L(s,\tau_i)=\prod_{i\geq2}L(s,\mu_i)$ 的任意 $u$ 阶极点，若不是 $L(s,\tau_1)$ 的极点，则为 $L(s,\pi)$ 的至少 $u$ 阶极点。5.5、5.6将证明这一点。

\paragraph*{5.5.}
现在采用3.5的记号。令 $n_1=p_1$、$n_2=n-n_1$、$\sigma_1=\pi(\mu_1)=\tau_1$、$\sigma_2=\pi(\mu_2\oplus\cdots\oplus\mu_r)$，以及 $\sigma=\sigma_1\times\sigma_2$。由归纳假设，
\[
 L(s,\sigma_2)=\prod_{2\leq i\leq r}L(s,\mu_i).\tag{5.5.1}
\]
与前面一样，$\pi$ 是 (3.5.3) 的商表示，其中 $P=MU$ 是型 $(n_1,n_2)$ 的抛物子群。$\pi$ 的矩阵系数由下列绝对收敛积分给出：
\[
 f(g)=\int_{M\backslash G}H(hg,h)\,dh
 =\int_{K\times\overline U}H(\bar u kg,k)\,d\bar k.\tag{5.5.2}
\]
这里 $\overline U=U_{\overline P}$、$\overline P={}^tP$\footnote{\textbf{译者注（本次校订）：}(5.5.2) 最右端的积分测度应为 $dk\,d\bar u$；正文的 $d\bar k$ 未定义。参见相同积分的 (3.3.8)。}，而 $H:G_n\times G_n\longrightarrow\C$ 是满足下列性质的任意函数：
\[
 H(u_1mg_1,\bar u_2mg_2)=H(g_1,g_2),\qquad
 u_1\in U,\ \bar u_2\in\overline U,\ m\in M;\tag{5.5.3}
\]
\[
 \text{对 }k_i\in K_n,\quad m\longmapsto H(mk_1,k_2)
 \text{ 是 }\sigma\otimes\delta_P^{1/2}\text{ 的矩阵系数};\tag{5.5.4}
\]
\[
 H\text{ 为 }C^\infty\text{ 函数，且对右作用是 }K_n\times K_n\text{ 有限的}.\tag{5.5.5}
\]
于是，若 $f$ 由(5.5.2)给出，则当 $\Re s$ 充分大时，$Z(\Phi,s+\tfrac12(n-1),f)$ 等于(3.5.5)（收敛问题留给读者，因为可用(3.6.6)、(3.6.7)处理）。反之，给定 $\Phi\in\mathcal S(n\times n,F)$ 及 $\sigma_i$ 的矩阵系数 $f_i$，令
\[
\begin{split}
 A(\Phi,s,f_1,f_2)={}&\int |\det m_1|^{s+(n_1-1)/2}|\det m_2|^{s+(n_2-1)/2}
 f_1(m_1)f_2(m_2)\\
 &\qquad\cdot\Phi\!\left[\begin{pmatrix}m_1&y\\xm_1&m_2+xy\end{pmatrix}\right]
 \,d^\times m_1\,d^\times m_2\,dx\,dy .
\end{split}\tag{5.5.6}
\]
当 $\Re s$ 充分大时，该重积分绝对收敛。我们将证明：
\begin{originaltheorem}{断言 5.5.7}
对任意 $\Phi$，商 $A(\Phi,s,f_1,f_2)/L(s,\pi)$ 可延拓为 $s$ 的整函数。若 $\Phi\in\mathcal S_0$，则它是 $s$ 的多项式。
\end{originaltheorem}
先设 $\Phi\in\mathcal S_0(n\times n,F)$。则 $K_n$ 上存在两个 $K$ 有限函数 $\xi,\xi'$，使(3.5.9)成立。令 $dh$ 为紧群 $K\cap P=K\cap\overline P=K\cap M$ 上的归一化 Haar 测度；对此群中的 $h$，写成
\[
 h=\begin{pmatrix}h_1&0\\0&h_2\end{pmatrix},\qquad h_i\in K_{n_i}.\tag{5.5.8}
\]
于是存在满足(5.5.3)--(5.5.5)的函数 $H$，使得
\[
\begin{split}
 H\!\left[\begin{pmatrix}m_1&0\\0&m_2\end{pmatrix}k',k\right]
 ={}&\delta_P^{1/2}\!\begin{pmatrix}m_1&0\\0&m_2\end{pmatrix}\\
 &\cdot\iint f_1(h_1m_1h_1'^{-1})f_2(h_2m_2h_2'^{-1})
 \xi'(h'k')\xi(hk)\,dh\,dh'.
\end{split}\tag{5.5.9}
\]
若 $f$ 是经(5.5.2)对应于 $H$ 的 $\pi$ 的矩阵系数，则与3.5中一样得到
\[
 A(\Phi,s,f_1,f_2)=Z(\Phi,s+\tfrac12(n-1),f).\tag{5.5.10}
\]
这对 $\Phi\in\mathcal S_0$ 证明了(5.5.7)。

还要指出，$A$ 也有类似(4.3.4)的关系，其中 $\Phi'$ 定义为
\[
 \Phi'\!\begin{pmatrix}m_1&x\\y&m_2\end{pmatrix}
 =\left.\frac{d}{dt}\Phi\!\begin{pmatrix}m_1e^{-t}&x\\ye^{-t}&m_2\end{pmatrix}\right|_{t=0},\tag{5.5.11}
\]
并且当 $\Phi\in\mathcal S_0$ 时，$\Phi'$ 也属于 $\mathcal S_0$。因此，$A$ 乘任意多项式所得仍是同型的积分，只须把 $\Phi$ 换为 $\Phi_1$。若 $\Phi\in\mathcal S_0$，则 $\Phi_1\in\mathcal S_0$。此外，映射 $\Phi\mapsto\Phi_1$ 连续。

现在我们将看到函数方程
\[
 B(\Phi^\wedge,1-s+\tfrac12(n-1),f_1^\vee,f_2^\vee)
 =\gamma(s,\pi,\psi)A(\Phi,s,f_1,f_2).\tag{5.5.12}
\]
这里 $B$ 由 $A$ 的定义中将 $(P,\sigma)$ 换成 $(\overline P,\widetilde\sigma)$ 得到。这样，显然可仿照(4.5.1)全纯性的证明来证明(5.5.7)（参见4.5的证明）。\footnote{译注：原书(5.5.12)左端的参数印有 $+\tfrac12(n-1)$；下页证明最后一式则把 $B$ 的该参数写为 $1-s$。此处分别保留扫描页所见形式，供读者对照。}

首先，显然可由 $(\overline P,\widetilde\sigma)$ 得到 $\widetilde\pi$，正如由 $(P,\sigma)$ 得到 $\pi$ 一样（参见3.3）。更准确地说，若 $f$ 是由(5.5.2)给出的 $\pi$ 的矩阵系数，则 $f^\vee$ 由(3.3.11)给出，而该处的 $\widetilde H$ 对 $(\overline P,\widetilde\sigma)$ 满足(5.5.3)--(5.5.5)。因此，每当 $f_i$ 是 $\widetilde\sigma_i$ 的矩阵系数时，令
\[
\begin{split}
 B(\Phi,s,f_1,f_2)={}&\int |\det m_1|^{s+(n_1-1)/2}|\det m_2|^{s+(n_2-1)/2}
 f_1(m_1)f_2(m_2)\\
 &\qquad\cdot\Phi\!\left[\begin{pmatrix}m_1+yx&ym_2\\x&m_2\end{pmatrix}\right]
 \,d^\times m_1\,d^\times m_2\,dx\,dy.
\end{split}\tag{5.5.13}
\]
于是，对 $\Phi\in\mathcal S_0$，(5.5.7)适用于 $B$，只是将 $L(s,\pi)$ 换为 $L(s,\widetilde\pi)$。现在给定 $\Phi\in\mathcal S_0$ 及 $\sigma_i$ 的矩阵系数 $f_i$；如上选取 $\xi,\xi',H$，使(5.5.9)成立，且(5.5.10)也成立，其中 $f$ 由(5.5.2)给出。则
\[
 \iint\Phi^\wedge[k^{-1}zk']\xi'(k)\xi(k')\,dk\,dk'=\Phi^\wedge(z).
\]
此外，按(3.3.11)，$\widetilde H$ 由 $\widetilde H(g_1,g_2)=H(g_2,g_1)$ 给出，故
\[
\begin{split}
 \widetilde H\!\left[\begin{pmatrix}m_1&0\\0&m_2\end{pmatrix}k',k\right]
 ={}&\delta_{\overline P}^{1/2}\!\begin{pmatrix}m_1&0\\0&m_2\end{pmatrix}\\
 &\cdot\iint f_1^\vee[h_1m_1h_1'^{-1}]f_2^\vee[h_2m_2h_2'^{-1}]
 \xi(h'k')\xi'(hk)\,dh\,dh'.
\end{split}
\]
由于 $f^\vee$ 由(3.3.11)给出，得到
\[
 B(\Phi^\wedge,1-s,f_1^\vee,f_2^\vee)
 =Z(\Phi^\wedge,1-s+\tfrac12(n-1),f^\vee).
\]
与(5.5.10)比较，即完成(5.5.12)及(5.5.7)的证明。

\paragraph*{5.6.}
令 $\mathcal A$ 为 $X=G_{n_1}\times M(n_2\times n_2,F)$ 上的函数空间，其中的函数关于第一变量具有紧支集（compact support），关于第二变量为 Schwartz 型函数。赋予 $\mathcal A$ 自然的拓扑。特别地，$C_c^\infty(X)$ 是 $\mathcal A$ 的稠密子集。另一方面，可将 $\mathcal A$ 视为 $M(n_1\times n_1,F)\times M(n_2\times n_2,F)$ 上 Schwartz 函数空间的子空间。对 $\Phi_1\in C_c^\infty(G_{n_1})$、$\Phi_{12}\in C_c^\infty(M(n_1\times n_2,F))$、$\Phi_{21}\in C_c^\infty(M(n_2\times n_1,F))$ 及 $\Phi_2\in\mathcal S(n_2\times n_2,F)$，可令
\[
 \Psi(m_1,x_2)=\Phi_1(m_1)\int\Phi_{21}(xm_1)\Phi_{12}(y)\Phi_2(x_2+xy)\,dx\,dy.\tag{5.6.1}
\]
令 $\mathcal A_0$ 为这些函数所张成的 $\mathcal A$ 的子空间；容易验证，$\mathcal A_0$ 包含 $C_c^\infty(X)$ 的一个稠密子空间，因此在 $\mathcal A$ 中稠密。

对 $\Psi\in\mathcal A$ 以及 $\sigma_1,\sigma_2$ 的矩阵系数 $f_1,f_2$，令
\[
\begin{split}
 U(\Psi,s,f_1,f_2)=\iint{}&\Psi(m_1,m_2)|\det m_1|^{s+(n_1-1)/2}
 |\det m_2|^{s+(n_2-1)/2}\\
 &\cdot f_1(m_1)f_2(m_2)\,d^\times m_1\,d^\times m_2.
\end{split}\tag{5.6.2}
\]
当 $\Re s$ 充分大时，此积分绝对收敛。由(5.5.7)可知：
\begin{originaltheorem}{断言 5.6.3}
对 $\Psi\in\mathcal A_0$，商 $U(\Psi,s,f_1,f_2)/L(s,\pi)$ 可延拓为 $s$ 的整函数。
\end{originaltheorem}
另一方面，我们将证明：
\begin{originaltheorem}{断言 5.6.4}
对 $\Psi\in\mathcal A$，商 $U(\Psi,s,f_1,f_2)/L(s,\sigma_2)$ 可延拓为 $s$ 的全纯函数。作为这样的函数，它是 $\mathcal A\times\C$ 上关于 $(\Psi,s)$ 的连续函数。
\end{originaltheorem}
注意，当 $\Psi$ 属于下列代数张量积时，(5.6.4)的第一项断言显然成立：
\[
 \mathcal A_1=C_c^\infty(G_n)\otimes\mathcal S(n_2\times n_2,F).\tag{5.6.5}
\]
因为若 $\Psi=\Psi_1\otimes\Psi_2$，则该商是一个绝对收敛积分与多项式 $Z(\Psi_2,s+\tfrac12(n_2-1),f_2)/L(s,\sigma_2)$ 的乘积。\footnote{译注：原书(5.6.5)的第一因子印作 $C_c^\infty(G_n)$，而本段函数的第一变量属于 $G_{n_1}$；第二因子印作 $\mathcal S$，未带下标 $0$。此处保留扫描页形式。}

再设 $s_0$ 为 $L(s,\sigma_2)$ 的 $u$ 阶极点；则存在 $f_1,f_2$ 及 $\Psi\in\mathcal A_1$，使(5.6.4)中的商在 $s_0$ 处有 $u$ 阶极点。\footnote{译注：原书此句写作“(5.6.4)中的商有 $u$ 阶极点”，与(5.6.4)所陈述的全纯性存在文字上的冲突；随后论证使用的是 $U(\Psi,s,f_1,f_2)$ 的极点。此处保留原句。}暂且承认(5.6.4)，假设 $s_0$ 不是 $L(s,\sigma_1)$ 的极点，且是 $L(s,\pi)$ 的一个阶数小于 $u$ 的极点。由于 $\mathcal A_0$ 在 $\mathcal A$ 中稠密，由(5.6.3)可知，亚纯依赖于 $s$ 的分布（distribution）$U(\,\cdot\,,s,f_1,f_2)$ 在 $s_0$ 处的极点阶数小于 $u$。标量亚纯函数 $U(\Psi,s,f_1,f_2)$ 也有同样性质。取 $\Psi\in\mathcal A_1$ 便得到矛盾。这证明了5.1。

因此只剩下证明(5.6.4)。如前所述，对 $\Psi\in\mathcal A_1$，第一项断言显然成立，此时该商是一个多项式与一个在竖直带中有界的函数的乘积。现在有关系
\[
 (as+b)U(\Phi,s,f_1,f_2)+U(\Phi',s,f_1,f_2)=0,\qquad a\ne0,\tag{5.6.6}
\]
其中当 $\Phi\in\mathcal A_1$ 时，$\Phi'\in\mathcal A_1$，且 $\Phi'$ 连续依赖于 $\Phi$。接着对充分小的 $\Re s$ 引入积分
\[
\begin{split}
 V(\Phi,s,f_1,f_2)=\int{}&\Psi(m_1,m_2)|\det m_1|^{s+(n_1-1)/2}
 |\det m_2|^{1-s+(n_2-1)/2}\\
 &\cdot f_1(m_1)f_2^\vee(m_2)\,d^\times m_1\,d^\times m_2.
\end{split}
\]
于是，对 $\Psi\in\mathcal A_1$，
\[
 V(\Psi^\wedge,s,f_1,f_2)=\gamma(s,\sigma_2,\psi)U(\Phi,s,f_1,f_2).
\]
$V$ 也满足类似(5.6.6)的关系。最后，$\mathcal A_1$ 在 $\mathcal A$ 中稠密。结合以上性质，沿用 (4.5.1) 的全纯性证明即可得到 (5.6.4)。这完成了5.1的证明。

\origsection{6}{整体理论（Global theory）}
现在域 $F$ 是一个 $A$ 域（$A$-field）。
\paragraph*{6.1.}
令 $\pi=\bigotimes_v\pi_v$ 为 $G_n(\A)$ 的不可约容许表示；与局部情形一样，这里的表示也按适当的代数或 $(\mathfrak g,K)$-模来理解，见 [I.P.S.]。于是 $\widetilde\pi=\bigotimes_v\widetilde\pi_v$。令
\[
 L(s,\pi)=\prod_vL(s,\pi_v),\qquad
 L(s,\widetilde\pi)=\prod_vL(s,\widetilde\pi_v).\tag{6.1.1}
\]
假设 $\pi$ 的中心拟特征 $\omega$ 在 $F^\times$ 上平凡。令
\[
 \varepsilon(s,\pi)=\prod_v\varepsilon(s,\pi_v,\psi_v).\tag{6.1.2}
\]
这里 $\psi=\prod\psi_v$ 是 $\A/F$ 上的非平凡特征；(6.1.2)中几乎所有因子均为 $1$，且乘积不依赖于 $\psi$。

\begin{originaltheorem}{定理 6.2（Theorem 6.2）}
假设 $\pi$ 是自守的。则无穷乘积(6.1.1)在某个右半平面内绝对收敛，并可延拓为整个复平面上 $s$ 的亚纯函数。延拓后满足
\[
 L(s,\pi)=\varepsilon(s,\pi)L(1-s,\widetilde\pi).
\]
若 $F$ 是数域，则它们仅有有限个极点，并在竖直带的无穷远处有界。若 $F$ 是函数域，且常数域有 $Q$ 个元素，则它们是 $Q^{-s}$ 的有理函数。
\end{originaltheorem}
\begin{proof}
若 $\pi$ 为尖点表示，当 $n=1$ 时，这是 [A.W. 2] 第 VII 章第6节的定理3、4、5；当 $n>1$ 时，这是 [R.G.-H.J.] 的定理13.8。若 $n>1$ 且 $\pi$ 不是尖点表示，则存在一个型 $(n_1,n_2,\ldots,n_r)$ 的标准抛物子群，以及对每个 $i$ 存在 $G_{n_i}$ 的不可约自守尖点表示 $\sigma_i$，满足下列条件：对任意位 $v$，$\pi_v$ 是诱导表示 $\xi_v=I(G_v,P_v;\sigma_{1v},\sigma_{2v},\ldots,\sigma_{rv})$ 的一个不可约成分（参见 [R.L.]）。因此，对任意位 $v$，
\[
\begin{split}
 L(s,\pi_v)&=P_v(s)\prod_iL(s,\sigma_{i,v}),\\
 L(s,\widetilde\pi_v)&=\widetilde P_v(s)\prod_iL(s,\widetilde\sigma_{i,v}),\\
 \gamma(s,\pi_v,\psi_v)&=\prod_i\gamma(s,\sigma_{i,v},\psi_v),
\end{split}
\]
其中当 $F_v=\R$ 或 $\C$ 时，$P_v$ 是 $s$ 的多项式；当 $F_v$ 非阿基米德且剩余域有 $q_v$ 个元素时，$P_v$ 是 $q_v^{-s}$ 的多项式。此外，对几乎所有 $v$，$\pi_v$ 及 $\sigma_{i,v}$ 均非分歧，因此由(3.6.5)，
\[
 P_v=\widetilde P_v=1,\qquad\varepsilon(s,\pi_v,\psi_v)=1
 \qquad\text{（几乎所有 }v\text{）}.
\]
从而
\[
\begin{split}
 L(s,\pi)&=\prod_vP_v(s)\prod_iL(s,\sigma_i),\\
 L(s,\widetilde\pi)&=\prod_v\widetilde P_v(s)\prod_iL(s,\widetilde\sigma_i),\\
 \varepsilon(s,\pi)&=\prod_v\frac{\widetilde P_v(1-s)}{P_v(s)}
 \prod_i\varepsilon(s,\sigma_i).
\end{split}
\]
由此即可得到各项断言；数域情形下 $L(s,\pi)$ 在竖直带中的有界性，再应用 Phragmén--Lindelöf 原理即可。\footnote{\textbf{译者注（本次校订）：}前一显示式中 $\varepsilon(s,\pi)$ 的多项式比值写反了。由 (2.7.5) 或直接代入 $\varepsilon=\gamma L/L(1-s,\widetilde\pi)$，应得 $\varepsilon(s,\pi)=\prod_v[P_v(s)/\widetilde P_v(1-s)]\prod_i\varepsilon(s,\sigma_i)$。这也正是 $L(s,\pi)=\varepsilon(s,\pi)L(1-s,\widetilde\pi)$ 所要求的比值。}
\end{proof}

\begin{thebibliography}{I.B.--A.Z. 3}
\bibitem[A.A.]{c23-AA} A. Andrianov, \emph{Zeta functions of simple algebras with non abelian characters}, Uspehi Mat. Nauk \textbf{23} (1968), 1--66 = Russian Math. Surveys \textbf{23} (4) (1968), 1--65.
\bibitem[I.B.--A.Z. 1]{c23-IBAZ1} I. N. Bernstein and A. V. Zelevinsky, \emph{Representations of the group $\GL(n,F)$ where $F$ is a non archimedean local field}, Uspehi Mat. Nauk \textbf{31} (3) (1976), 5--70 = Russian Math. Surveys \textbf{31} (3) (1976), 1--68.
\bibitem[I.B.--A.Z. 2]{c23-IBAZ2} I. N. Bernstein and A. V. Zelevinsky, \emph{Induced representations of the group $\GL(n)$ over a $p$-adic field}, Funkcional Anal. i Prilozen \textbf{10} (3) (1976), 74--75 = Functional Anal. Appl. \textbf{10} (3) (1976), 225--227.
\bibitem[I.B.--A.Z. 3]{c23-IBAZ3} I. N. Bernstein and A. V. Zelevinsky, \emph{Induced representations of reductive $p$-adic groups. II}, 预印本。
\bibitem[M.E.]{c23-ME} M. Eichler, \emph{Allgemeine Kongruenzklasseneinteilungen der Ideale einfacher Algebren über algebraischen Zahlkörpern und ihrer $L$-Reihen}, J. Reine Angew. Math. \textbf{179} (1938), 227--251.
\bibitem[G.F. 1]{c23-GF1} G. Fujisaki, \emph{On the zeta-functions of a simple algebra over the field of rational numbers}, Fac. Sci. Univ. Tokyo, Sect. 1 \textbf{7} (1958), 567--604.
\bibitem[G.F. 2]{c23-GF2} G. Fujisaki, \emph{On the $L$-functions of a simple algebra over the field of rational numbers}, Fac. Sci. Univ. Tokyo, Sect. 1 \textbf{9} (1962), 293--311.
\bibitem[S.G.]{c23-SG} S. Gelbart, \emph{Fourier analysis on matrix spaces}, Mem. Amer. Math. Soc., vol. 108, Providence, R.I., 1971.
\bibitem[R.G.]{c23-RG} R. Godement, \emph{Les fonction $\zeta$ des algèbres simples. I, II}, Sém. Bourbaki, 1957/1958, Exp. 171, 176.
\bibitem[R.G.--H.J.]{c23-RGHJ} R. Godement and H. Jacquet, \emph{Zeta functions of simple algebras}, Lecture Notes in Math., vol. 260, Springer-Verlag, 1972.
\bibitem[E.H.]{c23-EH} E. Hecke, \emph{Über Modulfunktionen und die Dirichletschen Reihen mit Eulerscher Produktenentwicklung. I, II}, Math. Ann. \textbf{114} (1937), 1--28; 316--351.
\bibitem[K.H.]{c23-KH} K. Hey, \emph{Analytische Zahlentheorie in Systemen hyperkomplexer Zahlen}, 学位论文，Hamburg, 1929.
\bibitem[K.I.]{c23-KI} K. Iwasawa, \emph{A note on $L$-functions}, Proc. Internat. Congr. Math., Cambridge, Mass., 1950, p. 322.
\bibitem[H.J. 1]{c23-HJ1} H. Jacquet, \emph{Zeta functions of simple algebras (local theory)}, Harmonic Analysis on Homogeneous Spaces, Mem. Amer. Math. Soc., Providence, R.I., 1973, pp. 381--386.
\bibitem[H.J. 2]{c23-HJ2} H. Jacquet, \emph{Generic representations}, Non-Commutative Harmonic Analysis, Lecture Notes in Math., vol. 587, Springer-Verlag, Berlin, 1977, pp. 91--101.
\bibitem[H.J.--R.L.]{c23-HJRL} H. Jacquet and R. Langlands, \emph{Automorphic forms on $\GL(2)$}, Lecture Notes in Math., vol. 114, Springer-Verlag, Berlin, 1970.
\bibitem[M.K.]{c23-MK} M. Kinoshita, \emph{On the $\xi$-functions of a total matrix algebra over the field of rational numbers}, J. Math. Soc. Japan \textbf{17} (1965), 374--408.
\bibitem[A.K.--G.Z.]{c23-AKGZ} A. W. Knapp and G. Zuckerman, \emph{Classification theorems for representations of semi-simple Lie groups}, Non-Commutative Harmonic Analysis, Lecture Notes in Math., vol. 587, Springer-Verlag, Berlin, 1977, pp. 138--159.
\bibitem[R.L.]{c23-RL} R. P. Langlands, \emph{On the classification of irreducible representations of real algebraic groups}, 油印讲义，Institute for Advanced Study, Princeton (1973).
\bibitem[H.M. 1]{c23-HM1} H. Maass, \emph{Über eine neue Art von nicht-analytischen automorphen Funktionen und die Bestimmung Dirichletscher Reihen durch Functionalgleichungen}, Math. Ann. \textbf{121} (1949), 141--183.
\bibitem[H.M. 2]{c23-HM2} H. Maass, \emph{Zetafunktionen mit Grössencharakteren und Kugelfunktionen}, Math. Ann. \textbf{134} (1957), 1--32.
\bibitem[G.M. 1]{c23-GM1} G. N. Maloletkin, \emph{Zeta functions of a semi-simple algebra over the field of rational numbers}, Dokl. Akad. Nauk SSSR \textbf{201} (3) (1971), 535--537 = Soviet Math. Dokl. \textbf{12} (6) (1971), 1678--1681.
\bibitem[G.M. 2]{c23-GM2} G. N. Maloletkin, \emph{Zeta functions of parabolic forms}, Mat. Sb. \textbf{86} (128) (1971), 622--643 = Math. USSR-Sb. \textbf{15} (4) (1971), 619--641.
\bibitem[A.S.]{c23-AS} A. Selberg, \emph{Harmonic analysis and discontinuous groups in weakly symmetric riemannian spaces with applications to Dirichlet series}, J. Indian Math. Soc. \textbf{20} (1956), 47--87.
\bibitem[H.S.]{c23-HS} H. Shimizu, \emph{On zeta functions of quaternion algebras}, Ann. of Math. (2) \textbf{81} (1965), 166--193.
\bibitem[G.S.]{c23-GS} G. Shimura, \emph{On Dirichlet series and abelian varieties attached to automorphic forms}, Ann. of Math. (2) \textbf{76} (1962).
\bibitem[E.S.]{c23-ES} E. Stein, \emph{Analysis in matrix spaces and some new representations of $\SL(N,\C)$}, Ann. of Math. (2) \textbf{86} (1967), 461--490.
\bibitem[T.T.]{c23-TT} T. Tamagawa, \emph{On the zeta functions of a division algebra}, Ann. of Math. (2) \textbf{77} (1963), 387--405.
\bibitem[J.T. 1]{c23-JT1} J. Tate, \emph{Fourier analysis in number fields and Hecke's zeta-functions}, Algebraic Number Theory, J. W. S. Cassels and A. Fröhlich 编，Thompson, 1967, pp. 305--347.
\bibitem[J.T. 2]{c23-JT2} J. Tate, \emph{Number theoretic background}, 本论文集第2册，pp. 3--26.
\bibitem[N.W.]{c23-NW} N. Wallach, \emph{Representations of reductive Lie groups}, 本论文集第1册，pp. 71--86.
\bibitem[G.W.]{c23-GW} G. Warner, \emph{Zeta functions on the real general linear group}, Pacific J. Math. \textbf{45} (1973), 681--691.
\bibitem[A.W. 1]{c23-AW1} A. Weil, \emph{Fonctions zeta et distributions}, Sém. Bourbaki, no. 312, (1965--66), 9 pp.
\bibitem[A.W. 2]{c23-AW2} A. Weil, \emph{Sur certains groupes d'opérateurs unitaires}, Acta Math. \textbf{111} (1964), 144--211.
\bibitem[A.W. 3]{c23-AW3} A. Weil, \emph{Basic number theory}, Springer-Verlag, Berlin, 1967.
\bibitem[A.W. 4]{c23-AW4} A. Weil, \emph{Dirichlet series and automorphic forms}, Lecture Notes in Math., vol. 189, Springer-Verlag, Berlin, 1971.
\bibitem[E.W.]{c23-EW} E. Witt, \emph{Riemann-Rochscher Satz und $\xi$-Funktion im Hyperkomplexen}, Math. Ann. \textbf{110} (1934), 12--28.
\bibitem[A.Z. 1]{c23-AZ1} A. V. Zelevinskii, \emph{Classification of irreducible non-cuspidal representations of the group $\GL(n)$ over $p$-adic fields}, Funkcional Anal. i Prilozen \textbf{11} (1) (1977), 67 = Functional Anal. Appl. \textbf{11} (1) (1977).
\bibitem[A.Z. 2]{c23-AZ2} A. V. Zelevinskii, \emph{Classification of representations of the group $\GL(n)$ over a $p$-adic field}, 预印本（俄文）。
\end{thebibliography}

\begin{flushright}
哥伦比亚大学（Columbia University）
\end{flushright}

